% concatenated from main.tex
\documentclass{article}

\usepackage[utf8]{inputenc}
\usepackage{hyperref}
\usepackage{enumitem}

\usepackage{array}


%\setlength{\oddsidemargin}{-18mm}\setlength{\evensidemargin}{-18mm}


%\setlength{\oddsidemargin}{-5mm} %\setlength{\evensidemargin}{-5mm}   %%
\setlength{\topmargin}{-20mm}
%\setlength{\textwidth}{194mm} \setlength{\textheight}{270mm}


\usepackage{tikz}
\usepackage{tikz-cd}
\usetikzlibrary{arrows}
\tikzcdset{arrow style=tikz,
diagrams={>={Stealth[round,length=4pt,width=6pt,inset=3pt]}}}

\usepackage{amsthm}

\usepackage{amsmath} 
\usepackage{amssymb}
%\usepackage{bbold}
\usepackage{amsfonts}
\usepackage{faktor}
\usepackage[T1]{fontenc}
\usepackage{faktor}
\usepackage{xfrac}
\usepackage{dsfont}
\usepackage{color}
 \usepackage{todonotes}
\usepackage[english]{babel}
\usepackage{mathtools}

\usepackage{mathdots}


 
\newtheorem{theorem}{Theorem}
\newtheorem{definition}[theorem]{Definition}
\newtheorem{conjecture}[theorem]{Conjecture}
\newtheorem{lemma}[theorem]{Lemma}
\newtheorem{notation}[theorem]{Notation}
\newtheorem{observation}[theorem]{Observation}
\newtheorem{proposition}[theorem]{Proposition}
\newtheorem{property}[theorem]{Property}
\newtheorem{example}[theorem]{Example}
\newtheorem{remark}[theorem]{Remark}
\newtheorem{corollary}[theorem]{Corollary}

\newtheorem*{theorem*}{Theorem}

\DeclareMathOperator{\ord}{ord}
\DeclareMathOperator{\rk}{rk}
\DeclareMathOperator{\length}{length}
\DeclareMathOperator{\Hess}{Hess}
\DeclareMathOperator{\syz}{syz}
\DeclareMathOperator{\chara}{char}
\DeclareMathOperator{\jet}{jet}
\DeclareMathOperator{\Ann}{Ann}
\DeclareMathOperator{\Spec}{Spec}
\DeclareMathOperator{\Bl}{Bl}
\DeclareMathOperator{\edim}{edim}
\DeclareMathOperator{\NZD}{NZD}
\DeclareMathOperator{\Frac}{Frac}
\DeclareMathOperator{\mult}{mult}
\DeclareMathOperator{\Hom}{Hom}
\DeclareMathOperator{\End}{End}
\DeclareMathOperator{\Tor}{Tor}


\newcommand{\parwy}{\partial^{\mathfrak{W}}_y}
\newcommand{\parwx}{\partial^{\mathfrak{W}}_x}
\newcommand{\muxyw}{\mu_{x,y}^{\mathfrak{W}}}



\numberwithin{theorem}{section}


\renewcommand{\mod}{\,\operatorname{mod}\,}




\title{The Milnor Number of One Dimensional Local Rings}
\author{Yotam Svoray}
\date{}

\AtEndDocument{\bigskip{\footnotesize%
  \textsc{ Department of Mathematics, University of Utah, UT 84112} \par  
  \textit{E-mail address}: \texttt{svoray@math.utah.edu} 
}}

\begin{document}

\maketitle

\begin{abstract}
    In this paper we present an analogue of the Milnor number for one dimensional local ring, and we show that it satisfies analogous properties to those of the Milnor number of plane curves over a field. In addition, we present two analogues of the semi-group of values of a one dimensional ring and show how they relate to our Milnor number. Finally, we use these tools and techniques to show we can relate the semigroups of one dimensional rings of finite Cohen-Macaulay type to those of the classical ADE singularities. 
\end{abstract}

  \tableofcontents


\section{Introduction}\label{sec:Intro3}

In this paper we study one dimensional Noetherian strictly Henselian reduced Nagata local ring in a manner similar to the study of algebraic curves over algebraically closed fields, and show that the have similar properties, behaviors, and classifications. We focus on two invariants and their analogues, the Milnor number of a curve and the semigroup of values of a curve.\\

Given a power series in two variables $f \in k[[x,y]]$ over an algebraically closed field $k$, the Milnor  number of $f$, denoted $\mu(f)$, is defined to be the $k-$dimension of the quotient ring $\frac{k[[x,y]]}{\langle \partial_x(f), \partial_y(f) \rangle}$. The Milnor number is an important invariant in the study of isolated singularities, defined by Milnor in~\cite{milnor1968singular}, which allows to understand their topology and algebraic behavior. For more information, see~\cite{greuel2007introduction, varchenko1985singularities}. Milnor also proved, using a topological argument, that in the case $k=\mathbb{C}$, we have an equality $\mu(f)=2\delta(f)-r(f)+1$, where $\delta(f)$ is the delta invariant of $f$ (defined to be the $k-$dimension of the quotient $\frac{\overline{\mathcal{O}_f}}{\mathcal{O}_f}$ with $\mathcal{O}_f=\frac{k[[x,y]]}{\langle f \rangle}$ and $\overline{\mathcal{O}_f}$ being its normalization) and $r(f)$ is the number of irreducible components, with an algebraic proof over arbitrary algebraically closed fields of characteristic zero given by Risler in~\cite{risler1971ideal}. Later,  Deligne in~\cite{deligne2006formule} proved that in the positive characteristic case, we have an inequality $\mu(f)\geq 2\delta(f)-r(f)+1$, and there has been much work on the study of the condition for which we have an equality in the positive characteristic case, see for example~\cite{melle2001pencils, nguyen2016invariants, boubakri2012invariants, greuel2012some, barroso2018milnor, barroso2016milnor}. In addition, in~\cite{barroso2022note}, the invariant $\overline{\mu}(f) :=2\delta(f)-r(f)+1$ was defined based upon Milnor's formula and studied. Therefore, based upon the invariant $\overline{\mu}$, in Section~\ref{sec:milnor} we define an analogous invariant for one dimensional local rings (with some technical conditions assumed) and show that it behaves in a similar fashion to the behavior of the Milnor number over an algebraically closed field of characteristic zero. In addition, we show how this Milnor number of a ring is much more independent of the characteristic of the residue field, compared to the traditional Milnor number (see, for example, discussion in the Introduction section of~\cite{boubakri2009hypersurface}). \\

In Section~\ref{sec:semigroup} we define two semigroup invariants, that are analogues of the classical semigroup of values of a curve and of a one dimensional local ring. The first, $\Gamma(R)$, which is a sub-semigroup of $\mathbb{N}$, is a generalization of the semigroup of values of a one dimensional domain, yet does not detect whether $R$ is a domain or not. The second, $\nu(R)$, is a sub-semigroup of $\mathbb{N}^{r(R)}$, where $r(R)$ is the number of minimal primes that $R$ has (which is the equivalent of "the number of irreducible components of a curve" for local rings), is more complicated to compute but encodes more information regarding $R$. We show how these semigroups are related to the Milnor invariant we defined and how we can use it to deduce information about the original ring. There has been much research on the relationship between one dimensional rings and algebraic curves with their semigroups of values, for example, see~\cite{assi2020numerical, barroso2012approach, barroso2022note, castellanos2005semigroup, campillo1994gorenstein, d2025value, de1987semigroup, kunz1970value, zariski2006moduli}. \\

Finally, in Section~\ref{sec:FCMT1} we use the Milnor number and the semigroups we define to relate between one dimensional Cohen-Macaulay local rings of finite Cohen-Macaulay type and the classical ADE singularities using a notion of "equisingularity", which is a generalized analogue of the "equisingularity of a curve" notion defined by Zariski in~\cite{zariski2006moduli}. There has been much research on one dimensional Cohen-Macaulay local rings of finite Cohen-Macaulay type. In particular, in this case, having finite Cohen-Macaulay type is equivalent to satisfying the Drodz-Ro\u{\i}ter conditions,  as first presented in~\cite{drozd1967commutative} and later expanded upon by \c{C}imen in~\cite{ccimen1998one}, based upon the work on Green and Reiner in~\cite{green1978integral}. Later, Greuel and Kr\"oning in~\cite{greuel1984einfache} classified one dimensional Cohen-Macaulay rings that contain an algebraically closed field  of characteristic zero as those containing and ADE singularity, a result that was later expanded upon by Wiegand in~\cite{wiegand1991curve, wiegand1994one} to one dimensional Cohen–Macaulay local ring that contain a perfect field of arbitrary characteristic. Baeth in~\cite{baeth2007krull} studied when such local rings satisfy the Krull-Schmidt theorem, based upon the work of Wiegand in~\cite{wiegand2023failure}. \\


We summarize the main results of this paper in the following theorem:

\begin{theorem*}
    \begin{enumerate}
        \item (Proposition~\ref{prop:uniqueness}) The Milnor number $\mu(R)$ (see Definition~\ref{def:Milnor}) is the unique invariant on the class of one dimensional Noetherian strictly Henselian reduced Nagata local ring that satisfies:
        \begin{enumerate}
        \item (Corollary~\ref{cor:reg}) $\mu(R)=0$ if and only if $R$ is a DVR, 
        \item (Lemma~\ref{lem:rho_delta}) If $r(R)=1$ (see Notation~\ref{notation:ringstuff}) then  we have $\mu(R)-\mu(R^{(1)})  = \rho(R)$ (see Definition~\ref{def:blowup}), 
        \item (Corollary~\ref{cor:lower_bound}) If $\mathfrak{p}_1, \dots, \mathfrak{p}_{r(R)}$ are the minimal primes of $R$ then we have $\mu(R) -1 = \sum_{i=1}^{r(R)} \left(\mu\left(\frac{R}{\mathfrak{p}_i}\right)-1\right) + 2\sum_{i <  j} i(\mathfrak{p}_i, \mathfrak{p}_j)$ (see Notation~\ref{notation:ringstuff}).
    \end{enumerate}
    \item (Proposition~\ref{prop:omp_lines}) $\mu(R) \geq (r(R)-1)^2$ with equality if and only if $R$ is an ordinary multiple point (see Definition~\ref{def:omp}). 
    \item (Theorem~\ref{thm:Morse}) If $R$ is analytically unramified, then  $\mu(R)=1$ if and only if $R$ is a double point (see Definition~\ref{def:doublept}). 
    \item (Proposition~\ref{prop:mu_vec}) If $R$ is Gorenstein then $\vec{\mu}(R)$ (see Definition~\ref{def:RelMIlnor}) is the minimal element in $\nu(R) \subset \mathbb{N}^{r(R)}$ (see Definition~\ref{def:nu}) such that $\vec{\mu}(R) + \mathbb{N}^{r(R)} \subset \nu(R)$.
    \item (Theorem~\ref{thm:ADE_curve}) If $R$ is Cohen-Macaulay and has finite Cohen-Macaulay type (see Definition~\ref{def:cm_type}), then $R$ is equisingular (see Definition~\ref{def:equidominated}) to an ADE type singularity (see Definition~\ref{def:ADE_curves}). 
    \end{enumerate}
\end{theorem*}




\noindent\textbf{Notations}: We denote $\mathbb{N}=\{0,1,2, \dots\}$ and the set of $r-$tuples over $\mathbb{N}$ by $\mathbb{N}^r$. We order $\mathbb{N}^r$ with the product order defined component-wise, that is, given two tuples $(\alpha_1, \dots, \alpha_r),(\beta_1, \dots, \beta_r) \in \mathbb{N}$ we say that $(\alpha_1, \dots, \alpha_r) \leq (\beta_1, \dots, \beta_r)$ if $\alpha_i \leq \beta_i$ for every $i$. Given some $\vec{a}_1, \dots, \vec{a}_n \in \mathbb{N}^r$, we denote by $\lceil \vec{a}_1, \dots, \vec{a}_n \rfloor$ the smallest (additive) semigroup of $\mathbb{N}^r$ that contains $\vec{a}_1, \dots, \vec{a}_n$. Given some $\vec{a} \in \mathbb{N}^r$, we denote $\vec{a} + \mathbb{N}^{r}$ the set of all $\vec{b} \in \mathbb{N}^r$ such that $\vec{b} \geq \vec{a}$. Similarly, we denote $\mathbb{N}_\infty = \mathbb{N} \cup \{\infty\}$ with its set of tuples $\mathbb{N}_\infty^r$ (together with the corresponding product order) and with sets of the form $\vec{a} + \mathbb{N}^r_\infty$.   \\

  

\noindent\textbf{Acknowledgments.} This work was done as part of the Author's PhD thesis under the guidance of Karl Schwede, and we wish to thank him for his guidance, help, and support. We wish to thank Tim Tribone for productive mathematical discussions and their inputs on some of the ideas presented in this paper. The author was partially supported by NSF grant DMS-2101800. 

\section{The Milnor Number $\mu(R)$}\label{sec:milnor}

In this section we define the Milnor number $\mu(R)$ of a one dimensional local ring and show that it behaves similarly to its analogue for algebraic curves. Before defining $\mu(R)$, we turn to some notions,  technical results, and remarks that we use throughout this paper. Most of these results are known to expertise, but we place them here for the sake of completeness and ease of reference.\\

Throughout this section, we assume that $(R, \mathfrak{m}, \kappa)$ is a one-dimensional Noetherian strictly Henselian reduced Nagata local ring with normalization $\overline{R}$. 

\begin{notation}\label{notation:ringstuff}
Let $(R, \mathfrak{m}, \kappa)$ be a one dimensional Noetherian local ring. Then: 
    \begin{enumerate}
        \item The delta invariant of $R$ is defined to be $\delta(R) = \length_R \left( \frac{\overline{R}}{R}\right)$, i.e. the length over $R$ of the quotient module $\frac{\overline{R}}{R}$. 
        \item Given a module $M$ over $R$, the multiplicity of $M$ over $R$ is defined to be $\mult_R(M)=\lim_{n \to \infty} \frac{\length_R\left(\frac{M}{\mathfrak{m}^n}\right)}{n}$, and the multiplicity of $R$ is defined to be $\mult(R)=\mult_R(R)$. 
        \item The embedding dimension of $R$ is defined to be $\edim(R) = \dim_{\kappa} \left( \frac{\mathfrak{m}}{\mathfrak{m}^2}\right)$.
        \item We denote by $\hat{R}$  the completion of $R$ (with respect to the maximal ideal $\mathfrak{m}$).
        \item We denote the number of minimal primes in $R$ (which is finite since $R$ is Noetherian) by $r(R)$.
        \item We denote by $\NZD(R)$ the set of non zero divisors of $R$.
        \item Given $I$ and $J$ be ideals in $R$. Then the intersection multiplicity of $I$ and $J$, denoted by $i(I, J)$, is the  the $R-$length of $\frac{R}{I+J}$. 
        \item The conductor of $R$ is the $R-$ideal $\mathfrak{c}_R=\Ann_R\left(\frac{\overline{R}}{R}\right)$ (which can in fact also be viewed as an $\overline{R}-$ideal)
    \end{enumerate}
\end{notation}

\begin{remark}\label{rem:techincal_ring}
\begin{enumerate}
    \item \textup{Note that since $R$ is one dimensional and $\kappa$ is infinite, then by~\cite{abhyankar1967local} we have that $\mult(R)\geq \edim(R)$.}
    \item \textup{Since $R$ is Nagata then it must be $J$-$1$, therefore we have that if $R$ is not regular then $R_\mathfrak{p}$ is regular for every $\mathfrak{p} \neq \mathfrak{m}$. This tells us that $R$ must have an isolated singularity.  }
    \item \textup{Since $R$ is reduced then we have that $\overline{R}$ is the product of the integral closure of $\frac{R}{\mathfrak{p}_i}$ over every minimal prime $\mathfrak{p}_1, \dots, \mathfrak{p}_{r(R)}$ of $R$.}
\end{enumerate}
\end{remark}



\begin{lemma}\label{lem:mult_sum_prim}
    If $\mathfrak{p}_1, \dots, \mathfrak{p}_{r(R)}$ are the minimal primes of $R$, then we have that $\mult(R)=\sum_{i=1}^r \mult\left(\frac{R}{\mathfrak{p}_i}\right)$. 
\end{lemma}

\begin{proof}
    Denote $\mathfrak{q}=\bigcap_{i=1}^{r(R)-1} \mathfrak{p}_i$. Then we that $\mathfrak{q}\cap \mathfrak{p}_{r(R)} = \{0\}$, and so we get a short exact sequence:
\begin{equation*}
    0 \to R \to \frac{R}{\mathfrak{q}} \oplus \frac{R}{\mathfrak{p}_{r(R)}} \to \frac{R}{\mathfrak{q}+\mathfrak{p}_{r(R)}} \to 0. 
\end{equation*}
Therefore, by applying multiplicity to the exact sequence (see Lemma 43.15.2 in~\cite{stacks-project}), we can conclude that $\mult(R) + \mult(\frac{R}{\mathfrak{p}_{r(R)}+\mathfrak{q}}) = \mult(\frac{R}{\mathfrak{q}} \oplus \frac{R}{\mathfrak{p}_{r(R)}})$. Yet, since $\frac{R}{\mathfrak{p}_{r(R)}+\mathfrak{q}}$ is zero dimensional, we get that  $\mult(\frac{R}{\mathfrak{p}_{r(R)}+\mathfrak{q}})=0$, and so the result follows by induction on $r(R)$ (with respect to the ring $\frac{R}{\mathfrak{q}}$). 
\end{proof}

\begin{corollary}\label{cor:mult_r}
    $\mult(R) \geq r(R)$.
\end{corollary}

\begin{proof}
    This follows directly from Lemma~\ref{lem:mult_sum_prim}, noting that $\mult \left(\frac{R}{\mathfrak{p}_i}\right) \geq 1$ for every $i$. 
\end{proof}

% \begin{lemma}\label{lem:ideal_edim}
% \begin{enumerate}
%     \item Let $I$ be a proper ideal of $R$. Then $\edim(R) \geq \edim\left(\frac{R}{I}\right)$ and we have an equality if $I \subset \mathfrak{m}^2$. 
%     \item  If $\mathfrak{p}_1, \dots, \mathfrak{p}_{r(R)}$ are the minimal primes of $R$ then $\edim(R) \leq \sum_i \edim \left(\frac{R}{\mathfrak{p}_i} \right)$.
% \end{enumerate}
% \end{lemma}


% \begin{proof}
%     For the first part, note that the maximal ideal of $\frac{R}{I}$ is $\overline{\mathfrak{m}}=\frac{\mathfrak{m}}{I}$, and so $\overline{\mathfrak{m}}^2 = \frac{\mathfrak{m}^2+I}{I}$. Therefore we can conclude that $\frac{\overline{\mathfrak{m}}}{\overline{\mathfrak{m}}^2} = \frac{\overline{\mathfrak{m}}}{\overline{\mathfrak{m}}^2 +I}$, and so we have a surjective map $\frac{{\mathfrak{m}}}{{\mathfrak{m}}^2} \to \frac{\overline{\mathfrak{m}}}{\overline{\mathfrak{m}}^2}$. By taking the dimension over $\kappa$ we can conclude that $\edim(R) \geq \edim(\frac{R}{I})$. If, in addition, $I \subset \mathfrak{m}^2$, then $I + \mathfrak{m}^2 = \mathfrak{m}^2$ and so we have an equality of the dimensions. For the second part,  the result follows. 
% \end{proof}


% \begin{lemma}\label{lem:milnor_finite}
%     If $R$ is not regular then $R_\mathfrak{p}$ is regular for every $\mathfrak{p} \neq \mathfrak{m}$.
% \end{lemma}

% \begin{proof}
% Since $R$ is Nagata then it must be $J$-$1$, and so the singular locus of $R$ (i.e. the set of all primes $\mathfrak{p} \subset R$ such that $R_\mathfrak{p}$ is not regular) is a closed subset $V(I) \subseteq \Spec (R)$. Since $R=R_\mathfrak{m}$ is not regular, then we have that $\mathfrak{m} \in V(I)$, and thus it is non empty. In addition, since $R$ is one dimensional and since the singular locus is not empty, then $\dim\left( \frac{R}{I} \right) \in \{0,1\}$. Yet, if $\dim(\frac{R}{I})=1$ then we can conclude that there exists some prime ideal $\mathfrak{p} \in V(I)$ such that $\dim(\frac{R}{\mathfrak{p}})=1$, and so $R_\mathfrak{p}$ must be non-reduced, and so we would have that $R$ is non reduced as well. 
% \end{proof}

\begin{lemma}\label{cor:edim_bound}
    $\delta(R) \geq \mult(R)-1$. 
\end{lemma}

\begin{proof}
    By applying length over $R$ to the short exact sequence of $R-$modules 
    \begin{equation*}
        0 \to \frac{R}{\mathfrak{m}} \to \frac{\overline{R}}{\mathfrak{m}} \to \frac{\overline{R}}{R} \to 0,
    \end{equation*}
    \noindent we get that $\delta(R)=\length_R\left(\frac{\overline{R}}{\mathfrak{m}}\right)-1$. Yet, we have that $\length_R\left(\frac{\overline{R}}{\mathfrak{m}}\right) \geq \length_R\left(\frac{\overline{R}}{\mathfrak{m}\overline{R}}\right)$ which equals to $\mult(R)$ (see Section 11.2. in~\cite{huneke2006integral}). 
\end{proof}

The following is a stronger version of the Cohen structure theorem, that we use extensively throughout this paper when classifying $R$ based upon its invariants: 

\begin{lemma}\label{lem:edim_CST}
    For every $n \geq \edim(R)$, there exists some complete regular local ring $(S, \mathfrak{n})$ with a surjective map $S \to \hat{R}$ such that $\dim(S)=n$ (where by $\dim$ we mean the Krull dimension of $S$).
\end{lemma}

\begin{proof}
    By looking at the ring of power series over a complete regular local ring, it is enough to prove the lemma in the case where $n=\edim(R)$. From the Cohen structure theorem (see, e.g. Theorem 10.160.8 in~\cite{stacks-project} or Corollary 28.3. in~\cite{matsumura1970commutative}), there exists some complete regular local ring $(S, \mathfrak{n})$ of some dimension $N\geq \edim(R)$ together with a surjective map $S \to \hat{R}$. Therefore, we have a surjective map of vector spaces $\frac{\mathfrak{n}}{\mathfrak{n}^2} \to \frac{\mathfrak{m}}{\mathfrak{m}^2}$ whose kernel is of dimension $N-\edim(R)$. Let $f_1, \dots, f_{N-\edim(R)}$ be elements in $S$ such that their image mod $\mathfrak{n}^2$ generate the kernel of this map of vector spaces, and set $S_1=\frac{S}{\langle f_1, \dots, f_{N-\edim(R)} \rangle}$. We have that $S_1$ is a complete ring of dimension $\edim(R)$ (by the Krull principal ideal theorem) and the map $S \to \hat{R}$ factors through $S_1$ giving us a surjective map $S_1 \to \hat{R}$. In addition, since the maximal ideal of $S_1$ is $\frac{\mathfrak{n}}{\langle f_1, \dots, f_{N-\edim(R)} \rangle}$, we have that $S_1$ is regular as its Zariski cotangent space is $\frac{\mathfrak{n}}{\langle f_1, \dots, f_{N-\edim(R)} \rangle + \mathfrak{n}^2}$, which is a vector space of dimension $\edim(R)$. 
\end{proof}

\begin{remark}
    \textup{Lemma~\ref{lem:edim_CST} gives us that $\hat{R}$ is isomorphic to some quotient $\frac{S}{I}$ for some complete regular local ring $S$. This observation is important when classifying rings $R$ (up to completion). }
\end{remark}

    We are now ready to define the Milnor number of $R$, inspired by the definition of $\overline{\mu}$ in~\cite{barroso2022note}:

\begin{definition}\label{def:Milnor}
 We define the Milnor number of $R$ to be $\mu(R)= 2 \delta(R)-r(R) + 1$.
\end{definition}

\begin{remark}\label{rem:conductor}
\begin{enumerate}
    \item \textup{By Corollary 12.2.4 in~\cite{huneke2006integral} and by Proposition 4.4.7 in~\cite{campillo2006algebroid} we have that $R$ is Gorenstein if and only if   $2\length_R(\frac{R}{\mathfrak{c}_R}) = \length_R(\frac{\overline{R}}{\mathfrak{c}_R}) $. In addition, since we have a short exact sequence
    \begin{equation*}
        0 \to \frac{R}{\mathfrak{c}_R} \to \frac{\overline{R}}{\mathfrak{c}_R} \to  \frac{\overline{R}}{R} \to 0, 
    \end{equation*}
    \noindent we can conclude that this is also equivalent to having $2 \delta(R)=  \length_R(\frac{\overline{R}}{\mathfrak{c}_R})$. }
    \item \textup{Note that $\mu(R)$ is finite since $\delta(R)$ is finite. This is true since $R$ is a one-dimensional reduced Nagata ring, and so $\overline{R}$ is a normal Nagata ring such that $R \to \overline{R}$ is a finite map. In particular, the support of $\frac{\overline{R}}{R}$ as an $R-$module is zero dimensional, and so it must have a finite length. }
    \item \textup{Remark~\ref{rem:techincal_ring} (together with the previous item) can be though of as an analogue of Lemma 2.3. in Chapter I of~\cite{greuel2007introduction}, i.e. that the Milnor number is finite if and only if our hypersurface has an isolated singularity. (Note that in our curve case, having an isolated singularity is equivalent to being reduced).  }
    \item \textup{Assuming the $R$ is analytically unramified, since $R$ is strictly Henselian, then from Lemma 33.39.5. and Lemma 33.39.6. of~\cite{stacks-project} we get that $\delta(R)=\delta(\hat{R})$ and $r(R)=r(\hat{R})$. Therefore we can conclude that $\mu(R)=\mu(\hat{R})$.  }
\end{enumerate}
\end{remark}

The following Lemma is an analogue of Lemma 33.40.5. and of  Lemma 33.40.6. in~\cite{stacks-project}:

\begin{lemma}\label{lem:milnor_nonneg}
    $\delta(R) \geq r(R)-1$ and so $\mu(R) \geq \delta(R) \geq 0$. 
\end{lemma}

\begin{proof}
     Let $\mathfrak{p}_1, \dots, \mathfrak{p}_{r(R)}$ be the minimal primes of $R$ and denote the residue field of $\frac{R}{\mathfrak{p}_i}$ by $\kappa_i$. Then we have a surjective map $\frac{\overline{R}}{R} \to \frac{\prod_{i=1}^{r(R)} \kappa_i}{\kappa}$. Note that since $\frac{R}{\mathfrak{p}_i}$ is a quotient ring of $R$ we must have that $\kappa \subset \kappa_i$, and so we can conclude that $\dim_\kappa \left(\frac{\prod_{i=1}^{r(R)}  \kappa_i}{\kappa}\right) \geq r(R)-1$. Therefore we have that $\delta(R) \geq r(R)-1$, and so $\mu(R) = \delta(R) + (\delta(R)-r(R)+1) \geq \delta(R)$. 
\end{proof}



We now turn to presenting an analogue of Hironaka's lemma (first proven by Hironaka for projective curves in~\cite{hironaka1957arithmetic}), which we use to classify $R$ for specific values of $\mu(R)$. For more information on Hironaka's Lemma, see Lemma 2.1 in~\cite{campillo1983hamburger}, Lemma 1.2.2 in~\cite{buchweitz1980milnor}, or Lemma 3.32. of Chapter I in~\cite{greuel2007introduction}. 



\begin{proposition}[Hironaka's Lemma] \label{thm:branch}
    Let $\mathfrak{p}_1, \dots, \mathfrak{p}_{r(R)}$ be the minimal primes of $R$. Then we have that 
    \begin{equation*}
        \delta(R) = \sum_{i=1}^{r(R)}  \delta\left(\frac{R}{\mathfrak{p}_i}\right) + \sum_{i <  j} i(\mathfrak{p}_i, \mathfrak{p}_j).
    \end{equation*}
    In addition, $\delta(R) \geq \frac{r(R)(r(R)-1)}{2}$. 
\end{proposition}
    

\begin{proof}
As in Lemma~\ref{lem:mult_sum_prim}, denote $\mathfrak{q}=\bigcap_{i=1}^{r(R)-1} \mathfrak{p}_i$. Then we that $\mathfrak{q}\cap \mathfrak{p}_{r(R)} = \{0\}$, and so we get a short exact sequence:
\begin{equation*}
    0 \to R \to \frac{R}{\mathfrak{q}} \oplus \frac{R}{\mathfrak{p}_{r(R)}} \to \frac{R}{\mathfrak{q}+\mathfrak{p}_{r(R)}} \to 0. 
\end{equation*}
\noindent Therefore, we can conclude that 
\begin{equation*}
    \delta(R)=\length_R\left(\frac{\overline{R}}{R}\right) = \length_R\left(\frac{\overline{R}}{\frac{R}{\mathfrak{q}} \oplus \frac{R}{\mathfrak{p}_{r(R)}} }\right) + \length_R \left(\frac{\frac{R}{\mathfrak{q}} \oplus \frac{R}{\mathfrak{p}_{r(R)}} }{R} \right).
\end{equation*}

\noindent Yet, the normalization of $\frac{R}{\mathfrak{q}} \oplus \frac{R}{\mathfrak{p}_{r(R)}}$ is exactly $\overline{R}$, that is, $\overline{R} = \overline{\left(\frac{R}{\mathfrak{q}}\right)} \times \overline{\left(\frac{R}{\mathfrak{p}_{r(R)}}\right)}$. In addition, we have that


\begin{equation*}
    \frac{\frac{R}{\mathfrak{q}} \oplus \frac{R}{\mathfrak{p}_{r(R)}} }{R}  = \frac{R}{\mathfrak{q} + \mathfrak{q}_{r(R)}}. 
\end{equation*}
\noindent Therefore, we can conclude that 

\begin{equation*}
    \delta(R) = \delta\left(\frac{R}{\mathfrak{q}}\right) + \delta\left(\frac{R}{\mathfrak{p}_{r(R)}} \right) + i_R(\mathfrak{p}_{r(R)}, \mathfrak{q}).
\end{equation*}
\noindent Since the minimal primes of $\frac{R}{\mathfrak{p}_{r(R)}}$ correspond to $\mathfrak{p}_1, \dots, \mathfrak{p}_{r(R)-1}$, we repeat this process for $\frac{R}{\mathfrak{p}_{r(R)}}$, and the result would follow. For the second part, observe that $i(\mathfrak{p}_i, \mathfrak{p}_j)>0$ for every $i  \neq j$, and so $\sum_{i <  j} i(\mathfrak{p}_i, \mathfrak{p}_j) \geq \frac{r(R)(r(R)-1)}{2}$. 
\end{proof}



From Proposition~\ref{thm:branch} we can conclude an analogous result for $\mu(R)$, which is an analogue of Property 2.3. in~\cite{ploski1995milnor}, of Corollary 1.2.3 in~\cite{buchweitz1980milnor}, and of item 2 of Proposition 1.1 in~\cite{barroso2018milnor}:

\begin{corollary}\label{cor:lower_bound}
    Let $\mathfrak{p}_1, \dots, \mathfrak{p}_{r(R)}$ be the minimal primes of $R$. Then 
    \begin{equation*}
        \mu(R) -1 = \sum_{i=1}^{r(R)} \left(\mu\left(\frac{R}{\mathfrak{p}_i}\right)-1\right) + 2\sum_{i <  j} i(\mathfrak{p}_i, \mathfrak{p}_j).
    \end{equation*}
    In particular, $\mu(R) \geq (r(R)-1)^2$. 
\end{corollary}

\begin{proof}
    The first part follows directly from Proposition~\ref{thm:branch} together with the definition of $\mu(R)=2\delta(R)-r(R)+1$. For the second part, as in Proposition~\ref{thm:branch}, we have that $2\sum_{i <  j} i(\mathfrak{p}_i, \mathfrak{p}_j) \geq r(R)(r(R)-1)$ and so $\mu(R) \geq r(R)(r(R)-1) - r(R) +1 = (r(R)-1)^2$.
\end{proof}


\begin{remark}
    \textup{Note that Corollary~\ref{cor:lower_bound} provides us with an alternative proof to Lemma~\ref{lem:milnor_nonneg} (i.e. that $\mu(R)$ is non-negative, as $r(R)\geq 1$). }
\end{remark}

We can use Corollary~\ref{cor:lower_bound} to provide a classification of DVR as rings that satisfy $\mu(R)=0$, which is an analogue of Lemma 2.44. in Chapter I of~\cite{greuel2007introduction} and of item 5 of Proposition 1.1 in~\cite{barroso2018milnor}. 

\begin{corollary}\label{cor:reg}
    $\mu(R)=0$ if and only if $R$ is a DVR. 
\end{corollary}
    

\begin{proof}
    By Corollary~\ref{cor:lower_bound}, if $\mu(R)=0$, we must get that $r(R)=1$ (as $r(R)>0$), and so we can conclude that $0=\mu(R)=\delta(R)$. Therefore $R$ must normal, and a normal one dimensional local ring must be a DVR. 
\end{proof}

The following theorem is an analogue version of Morse's lemma, as presented in Theorem 2.46 in~\cite{greuel2007introduction} (which is based upon~\cite{morse1934calculus, varchenko1985singularities}). Note that it is also a characteristic-free one dimensional analogue of Proposition 3.3. in~\cite{carvajal2019covers}, together with its generalization in Proposition 3.14 of~\cite{svoray2025ade}. 

\begin{definition}\label{def:doublept}
    We say that $(R , \mathfrak{m})$ is a double point if there exists some $2-$dimensional complete regular local ring $(S, \mathfrak{n})$ such that $\mathfrak{n} = \langle x,y \rangle$ and $\hat{R}$ is isomorphic to $\frac{S}{\langle xy \rangle}$. 
\end{definition}


\begin{theorem}\label{thm:Morse}
    Assume that $R$ is analytically unramified. Then $\mu(R)=1$ if and only if $R$ is a double point. 
\end{theorem}

\begin{proof}
    If $\mu(R)=1$ then by Corollary~\ref{cor:lower_bound} we have that $1 \geq (r(R)-1)^2$, and so $r(R) \in \{1,2\}$. Yet, $r(R)$ can not be $1$ since as $\mu(R)=2\delta(R)-r(R)+1$ then we would conclude that $1 = \mu(R) = 2\delta(R)$, which is impossible as $\delta(R)$ is an integer. Therefore we must have that $r(R)=2$. Denote the minimal primes of $R$ by $\mathfrak{p}$ and $\mathfrak{q}$. This gives us that $1=\mu(R)=2\delta(R)-1$, and so $\delta(R)=1$. Thus by Proposition~\ref{thm:branch} we can conclude that 
    \begin{equation*}
        1=\delta(R)=\delta\left(\frac{R}{\mathfrak{p}}\right) + \delta\left(\frac{R}{\mathfrak{q}}\right) + i(\mathfrak{p}, \mathfrak{q}).
    \end{equation*}

    \noindent Note that $i(\mathfrak{p}, \mathfrak{q}) >0$, as otherwise we would get that $R=\mathfrak{p}+\mathfrak{q}$ which is impossible. Therefore we must get that $\delta\left(\frac{R}{\mathfrak{p}}\right) =\delta\left(\frac{R}{\mathfrak{q}}\right)=0$ and $i(\mathfrak{p}, \mathfrak{q}) =1$. So, from Corollary~\ref{cor:reg}  we can conclude that $\frac{R}{\mathfrak{p}}$ and $\frac{R}{\mathfrak{q}}$ are DVR's with $i(\mathfrak{p}, \mathfrak{q})=1$, which is equivalent to $\mathfrak{p}+\mathfrak{q}=\mathfrak{m}$. Thus, we have that $\mult\left(\frac{R}{\mathfrak{p}}\right) = \mult\left(\frac{R}{\mathfrak{q}}\right)=1$, and so by Lemma~\ref{lem:mult_sum_prim} we can conclude that $\mult(R)=\mult\left(\frac{R}{\mathfrak{p}}\right) + \mult\left(\frac{R}{\mathfrak{q}}\right) = 2$. Yet, by Remark~\ref{rem:techincal_ring}, since $\edim(R) \leq \mult(R)=2$ and since $R$ is not a DVR (and therefore $\edim(R)>1$), we must have that $\edim(R)=2$. \\

    Therefore, by Lemma~\ref{lem:edim_CST} we have that $\hat{R} \cong \frac{S}{I}$ for some two dimensional complete regular ring $(S, \mathfrak{n})$ and some ideal $I \subset S$. Since $R$ is analytically unramified, then $\hat{R}$ is reduced. As $R$ has two minimal primes, then by item 4 in Remark~\ref{rem:conductor}, there are two prime ideals $\mathfrak{p}_1, \mathfrak{q}_1 \subset S$ such that $I = \sqrt{I} = \mathfrak{p}_1 \cap \mathfrak{q}_1$. Since $S$ is two dimensional and regular we can conclude that $\mathfrak{p}_1 =  \langle s_1 \rangle  $ and $\mathfrak{q}_1 =  \langle s_2 \rangle  $ for some $s_1, s_2 \in \mathfrak{n} \setminus \mathfrak{n}^2$. Since $r(R)=2$ we must have that $\mathfrak{p}_1 \neq \mathfrak{q}_1$, therefore $\langle s_1, s_2 \rangle = \mathfrak{n}$ and we can conclude that $I = \langle s_1 \rangle \cap \langle s_2 \rangle = \langle s_1s_2 \rangle$. 
\end{proof}

\begin{remark}
    \textup{From Corollary $2$ in~\cite{ooishi1991conductor} together with Theorem~\ref{thm:Morse} and item 1 in Remark~\ref{rem:conductor} we can conclude that if $R$ is analytically unramified and Gorenstein then $\mult(R)=\edim(R)$ if and only if $R$ is either a double point or a DVR. }
\end{remark}

We can use Theorem~\ref{thm:Morse} to conclude for which rings $R$ we have equality in Lemma~\ref{lem:milnor_nonneg}: 

\begin{corollary}
    Assume that $R$ is analytically unramified. Then the following are equivalent:
    \begin{enumerate}
        \item $\mu(R)=\delta(R)$.
        \item $\delta(R)=r(R)-1$. 
        \item $R$ is either a DVR or a double point. 
    \end{enumerate}
\end{corollary}

\begin{proof}
    The equivalence of $1$ and $2$ follows from Lemma~\ref{lem:milnor_nonneg}. Now, if $\delta(R)=r(R)-1$ then by Proposition~\ref{thm:branch} we can conclude that $r(R)-1 = \delta(R) \geq \frac{r(R)(r(R)-1)}{2}$, and so either $r(R)=1$ or $r(R)=2$. If $r(R)=1$ then $\delta(R)=0$ and so $R$ is a DVR (as in Corollary~\ref{cor:reg}), and if $r(R)=2$ then $\delta(R)=1$ and so we can conclude that $\mu(R)=1$, which gives us that $R$ is a double point by  Theorem~\ref{thm:Morse}.
\end{proof}

Inspired by Theorem~\ref{thm:Morse} we now turn to looking at two generalizations of the double point, namely the ordinary multiple point and the $A_k$ singularities. We show that they have similar classification properties based on the values of $\mu(R)$ and $\delta(R)$.  

\begin{definition}\label{def:omp}
    We say that $R$ is an ordinary multiple point if it satisfies $\mu(R)=(r(R)-1)^2$.  
\end{definition}

\begin{remark}\label{rem:omp}
\begin{enumerate}
    \item \textup{Note that by the definition of $\mu(R)$ we have that $R$ is an ordinary multiple point if and only if $\delta(R)=\frac{r(R)(r(R)-1)}{2}$. In addition, if we denote the minimal primes of $R$ by $\mathfrak{p}_1, \dots, \mathfrak{p}_{r(R)}$, then by Proposition~\ref{thm:branch} and Corollary~\ref{cor:lower_bound} together with Corollary~\ref{cor:reg} we can conclude that $R$ is an ordinary multiple point if and only if $\frac{R}{\mathfrak{p}_i}$ is a DVR for every $i$ and that $i(\mathfrak{p}_i, \mathfrak{p}_j)=1$ for every $i \neq j$. }
    \item \textup{Notice that an ordinary multiple point can be though of as a ring with the minimal value for $\mu(R)$, i.e. the equality cases in Proposition~\ref{thm:branch} and Corollary~\ref{cor:lower_bound}. }
\end{enumerate}
\end{remark}

The following proposition is an analogue of Lemma 1.2.4 in~\cite{buchweitz1980milnor}:

\begin{proposition}\label{prop:omp_lines}
If $R$ is an ordinary multiply point then $\edim(R) \leq r(R)$. In addition, the following are equivalent:
\begin{enumerate}
    \item $R$ is an ordinary multiple point with $r(R)=\edim(R)$, 
    \item There exist some regular sequence $x_1, \dots, x_{r(R)} \in R$ that generate $\mathfrak{m}$ such that for every $i$ we have that $\mathfrak{p}_i = \langle x_1, \dots, x_{i-1}, x_{i+1}, \dots, x_{r(R)} \rangle$.
    \item There exists some $r(R)-$dimensional complete regular local ring $(S, \mathfrak{n})$ with $\mathfrak{n} =\langle x, \dots, x_{r(R)} \rangle$ such that $\hat{R}$ is isomorphic to $\frac{S}{\langle x_ix_j \colon i \neq j \rangle}$.
\end{enumerate}
\end{proposition}

%Then $\edim(R) \leq r(R)$. If, in addition, $R$ is analytically unramified, then we have equality if and only if .

\begin{proof}
    Denote the minimal primes of $R$ by $\mathfrak{p}_1, \dots, \mathfrak{p}_{r(R)}$. Then by Remark~\ref{rem:omp} we have that $\frac{R}{\mathfrak{p}_i}$ is a DVR for every $i$. since $\mathfrak{p}_1, \dots \mathfrak{p}_r$ are the minimal primes of $R$, we have that $\mathfrak{p}_1 \cap \dots \cap \mathfrak{p}_r = \langle 0 \rangle \subset \mathfrak{m}^2$. Therefore if we can look at the map of $\kappa-$vector spaces 
    \begin{equation*}
        \frac{{\mathfrak{m}}}{{\mathfrak{m}}^2} \to\bigoplus_{i=1}^{r(R)} \frac{\mathfrak{m}}{\mathfrak{p}_i + \mathfrak{m}^2}, 
    \end{equation*}
    \noindent defined by $x + \mathfrak{m}^2 \mapsto (x+ \mathfrak{p}_i + \mathfrak{m}^2)_{i=1}^{r(R)}$, then we claim this map is injective. If for every $i$ we have that $x+ \mathfrak{p}_i + \mathfrak{m}^2 = 0+ \mathfrak{p}_i + \mathfrak{m}^2$ then $x \in \cap_{i=1}^{r(R)} (\mathfrak{p}_i + \mathfrak{m}^2)$ and so for every $i$ we can find some $a_i \in \mathfrak{p}_i$ such that $x-a_i \in \mathfrak{m}^2$. Thus we have that $a_i-a_j \in \mathfrak{m}^2$, which gives us that $x \in \mathfrak{m}^2+ \cap_{i=1}^{r(R)} \mathfrak{p}_i$. As $\mathfrak{p}_1 \cap \cdots \cap \mathfrak{p}_{r(R)}=\langle 0 \rangle$, then $\mathfrak{p}_1 \cap \cdots \cap \mathfrak{p}_{r(R)} = 0 \mod \mathfrak{m}^2$, which gives us that $x \in \mathfrak{m}^2$. Thus, by taking the dimension over $\kappa$ to both sides, we can conclude that $\edim(R) \leq \sum_{i=1}^{r(R)} \edim \left( \frac{R}{\mathfrak{p}_i} \right) = r(R)$.\\
    
    Now, if $r(R)=\edim(R)$, we have that this map is an isomorphism. Therefore,  we can find some $x_1, \dots, x_{r(R)} \in \mathfrak{m} \setminus \mathfrak{m}^2$ that generate $\mathfrak{m}$ such that $x_i + (\mathfrak{p}_i +\mathfrak{m}^2)$ spans $\frac{\mathfrak{m}}{\mathfrak{p}_i + \mathfrak{m}^2}$ as a $\kappa-$vector space. Since for every $i \neq j$ we have that the image of $x_j$ in the vector $\frac{\mathfrak{m}}{\mathfrak{p}_i + \mathfrak{m}^2}$ is zero, we  must have that $x_j \in \mathfrak{p}_i$. Therefore, since $\frac{R}{\mathfrak{p}_i}$ is a DVR we can conclude that $\mathfrak{p}_i = \langle x_1, \dots, x_{i-1}, x_{i+1}, \dots, x_n\rangle$ for every $i$. 
\end{proof}

\begin{remark}\label{rem:omp_lines_edim}
    \textup{Proposition~\ref{prop:omp_lines} tells us that intuitively, an ordinary multiple point can be geometrically though of as "a collection of $r(R)$ lines space all intersecting at a single point".  }
\end{remark}
%he following Corollary and Proposition give us an analogue of Theorem~\ref{thm:Morse} for ordinary multiple points in the two extreme values of $\edim(R)$ (as described in Remark~\ref{rem:omp_lines_edim}): 

% \begin{corollary}
%     If $r(R)=\edim(R)$ then $R$ is an ordinary multiple point if and only if 
% \end{corollary}

% \begin{proof}
%     This follows directly from Proposition~\ref{prop:omp_lines} together with the second item of Remark~\ref{rem:techincal_ring}. 
% \end{proof}


% In particular, we have an equality $\edim(R)=r(R)$ if and only if we have an isomorphism of $\kappa-$vector spaces $\frac{{\mathfrak{m}}}{{\mathfrak{m}}^2} \cong \bigoplus_{i=1}^{r(R)} \frac{\mathfrak{m}}{\mathfrak{p}_i + \mathfrak{m}^2}$. 

%Now, if $\edim(R)=r(R)$ then by Remark~\ref{rem:techincal_ring} there exist some complete regular local ring $(S, \mathfrak{n})$ of dimension $r(R)$ such that $\hat{R} \cong \frac{S}{I}$. Since $R$ is analytically unramified, then $I$ is radical and so we can write $I = \cap_{i=1}^{r(R)} \mathfrak{q}_i$. By the previous part, we have an isomorphism of $\kappa-$vector spaces $\frac{{\mathfrak{n}}}{{\mathfrak{n}}^2} \cong \bigoplus_{i=1}^{r(R)} \frac{\mathfrak{n}}{\mathfrak{q}_i + \mathfrak{n}^2}$, therefore we can find a $\kappa-$basis $x_1+\mathfrak{n}^2, \dots, x_{r(R)} + \mathfrak{n}^2$ for $\frac{{\mathfrak{n}}}{{\mathfrak{n}}^2}$ such that $x_i + \mathfrak{n}^2$ span $\frac{\mathfrak{n}}{\mathfrak{q}_i + \mathfrak{n}^2}$ over $\kappa$. Therefore we can conclude that $x_1, \dots, x_n$ generate $\mathfrak{n}$ with $\mathfrak{q}_i = \langle x_1, \dots, x_{i-1}, x_{i+1}, \dots, x_n\rangle$ for every $i$, and the result follows.


% \begin{remark}\label{rem:omp}
%     \textup{Note that in Definition~\ref{def:omp} we have an equality of ideals $\langle x_ix_j \colon i \neq j \rangle = \bigcap_{i=1}^n \langle x_1, \dots x_{i-1}, x_{i+1}, \dots , x_n\rangle$. Therefore we can geometrically view an ordinary multiple point as . Yet, by Lemma~\ref{lem:ideal_edim} we have that the embedding dimension of an ordinary multiple point is exactly $r(R)$, and so an ordinary multiple point can not be written as a quotient of a two dimensional complete regular local ring. }
% \end{remark}



% \begin{proposition}\label{lem:ineq_edim_e}
% Assume that $R$ is analytically unramified. Then the following are equivalent:
%     \begin{enumerate}
%         \item  $\mu(R) = (r(R)-1)^2$,
%         \item , 
%         \item $R$ is an ordinary multiple point.
%     \end{enumerate}
% \end{proposition}

% \begin{proof}
%     Note that the equivalence of $1$ and $2$ follows directly from the definition of $\mu(R)$. We prove their equivalence to $3$ using induction on $r=r(R)$. Note that the case $r(R)=1$ follows from Corollary~\ref{cor:reg} and the case $r(R)=2$ follows from Theorem~\ref{thm:Morse}. 
    
%     . 
    
%     . Therefore we must have that $\edim\left(\frac{R}{\mathfrak{p}_i}\right)=1$ and so by Lemma~\ref{lem:ideal_edim} we can conclude that $\edim(R) \leq r(R)$. Thus, by Remark~\ref{rem:techincal_ring}, there exists some complete regular local ring $(S, \mathfrak{n})$ of dimension $r(R)$ and some ideal $I \subset \mathfrak{n}$ such that $\hat{R} \cong \frac{S}{I}$. In addition, since $R$ is analytically unramified we can conclude that $I$ must be a radical ideal. So $I = \bigcap_{i=1}^{r(R)} \mathfrak{q}_i$ where $\mathfrak{q}_1, \dots, \mathfrak{q}_{r(R)}$ are the minimal primes of $\frac{S}{I}$ (that correspond with $\mathfrak{p}_1, \dots, \mathfrak{p}_r$). Thus, if we denote $\mathfrak{p}=\bigcap_{i<n} \mathfrak{p}_i$, then by the proof of Proposition~\ref{thm:branch} we can conclude that $i(\mathfrak{p}_{r(R)}, \mathfrak{p})=\sum_{i<r(R)} i(\mathfrak{p}_{r(R)}, \mathfrak{p}_i)=r(R)-1$. Therefore we get that 
%     \begin{equation*}
%        \frac{r(r-1)}{2}=\delta(R)=\delta\left(\frac{S}{\mathfrak{p}_n}\right)+\delta\left(\frac{R}{\mathfrak{p}}\right)+i(\mathfrak{p}_n, \mathfrak{p})=0+\delta\left(\frac{R}{\mathfrak{p}} \right) + (r-1),
%     \end{equation*}
%     which gives us that $\delta\left(\frac{R}{\mathfrak{p}} \right)=\frac{(r-1)(r-2)}{2}$, and so by the induction step we can conclude that $\frac{R}{\mathfrak{p}}$ is an ordinary multiple point as well with $(r(R)-1)$ minimal primes.  Therefore, there exists some complete regular ring $(S_1, \mathfrak{n}_1)$ of dimension $r-1$ such that $\mathfrak{n}_1 = \langle y_1, \dots, y_{r-1} \rangle$ and the completion of $\frac{R}{\mathfrak{q}}$ is isomorphic to $\frac{S_1}{ \langle y_iy_j \rangle_{i \neq j}}$. But the completion of $\frac{R}{\mathfrak{p}}$ is isomorphic to $\frac{S}{\mathfrak{q}}$ where $\mathfrak{q}=\bigcap_{i<n} \mathfrak{q}_i$. Therefore by lifting $y_1, \dots, y_{r-1}$ via this isomorphism we can find a regular sequence $x_1, \dots, x_r$ of $S$ such that $\mathfrak{q}_i= \langle x_1, \dots x_{i-1}, x_{i+1}, \dots , x_n\rangle$ for every $i$, and by recalling Remark~\ref{rem:omp}, the result follows. 
% \end{proof}

    % Therefore, if we denote $\mathfrak{q}=\bigcap_{i<n} \mathfrak{p}_i$, then by the proof of Proposition~\ref{thm:branch} we can conclude that $i(\mathfrak{p}_r, \mathfrak{q})=\sum_{i<r} i(\mathfrak{p}_r, \mathfrak{p}_i)=r-1$. Therefore we get that 
    % \begin{equation*}
    %    \frac{r(r-1)}{2}=\delta(R)=\delta\left(\frac{R}{\mathfrak{p}_n}\right)+\delta\left(\frac{R}{q}\right)+i(\mathfrak{p}_n, \mathfrak{q})=0+\delta\left(\frac{R}{\mathfrak{q}} \right) + (r-1),
    % \end{equation*}
    % which gives us that $\delta\left(\frac{R}{\mathfrak{q}} \right)=\frac{(r-1)(r-2)}{2}$, and so by the induction step we can conclude that $\frac{R}{\mathfrak{q}}$ is an ordinary multiple point as well with $r_{\frac{R}{\mathfrak{q}}} = r(R)-1$.  Therefore, there exists some complete regular ring $(S_1, \mathfrak{n}_1)$ of dimension $r-1$ such that $\mathfrak{n}_1 = \langle y_1, \dots, y_{r-1} \rangle$ and the completion of $\frac{R}{\mathfrak{q}}$ is isomorphic to $\frac{S_1}{ \langle y_iy_j \rangle_{i \neq j}}$. Therefore $\mult\left(\frac{R}{\mathfrak{q}}\right)=r(R)-1$, and since $\frac{R}{\mathfrak{p}}$ is a DVR, by Lemma~\ref{lem:mult_sum_prim} we can conclude that $\mult(R)=r(R)$.\\

    % ....... Therefore, there exist some $r(R)$ dimensional complete regular local ring $(S, \mathfrak{n})$ such that $\frac{\hat{R}}{\mathfrak{p}}$, $\frac{\hat{R}}{\mathfrak{q}}$, and $\hat{R}$ are all quotient rings of $S$. Since $\frac{\hat{R}}{\mathfrak{q}}  \cong \frac{S_1}{ \langle y_iy_j \rangle_{i \neq j}}$ and since $\frac{R}{\mathfrak{p}}$ is a DVR,  we can find some $x, \dots, x_{r(R)}$ that generate $\mathfrak{n}$ such that $\frac{\hat{R}}{\mathfrak{p}} \cong \frac{S}{\langle x, \dots, x_{r(R)} \rangle}$ and $\frac{\hat{R}}{\mathfrak{q}} = \frac{S}{\langle x_{r(R)} \rangle \cdot \langle x_ix_j \rangle_{i \neq j<r(R)}}$. Thus we can conclude that $\hat{R} \cong \frac{S}{\langle x_ix_j \rangle_{i \neq j}}$ and the result follows.

% \begin{remark}\label{rem:mult_r}
%     \textup{By the proof of Proposition~\ref{lem:ineq_edim_e} we can conclude that we have an equality in Corollary~\ref{cor:mult_r}, i.e. we have $\mult(R)=r(R)$, if and only if $R$ is an ordinary multiple point. This is true since $\mult(R)=r(R)$ if and only if $\frac{R}{\mathfrak{p}_i}$ is a DVR for every $i$ and so.............}
% \end{remark}

\begin{proposition}\label{prop:omp_edim2}
    Assume $R$ is an ordinary multiple point. Then $\edim(R)=2$ if and only if $\hat{R}$ is isomorphic to the quotient ring $\frac{S}{\langle f\rangle}$ for some $2$-dimensional complete regular local ring $(S, \mathfrak{n})$ with $\mathfrak{n} = \langle x,y \rangle$ and where $f=\prod_{i=1}^{r(R)} (x+yu_i)$ for distinct unit $u_1, \dots, u_{r(R)} \in S$ such taht $u_i-u_j$ is a unit for every $i \neq j$.  
\end{proposition}


\begin{proof}
   We can assume that $r(R) > 2$ the case $r(R)=1$ follows from Corollary~\ref{cor:reg} and the case $r(R)=2$ follows from Theorem~\ref{thm:Morse}. As $\edim(R)=2$ then by Lemma~\ref{lem:edim_CST} we have that $\hat{R}$ is isomorphic to some quotient $\frac{S}{\langle f \rangle}$ where $(S,\mathfrak{n})$ is a two dimensional complete regular local ring and $f \in \mathfrak{n}$. By Remark~\ref{rem:omp} we have that $\frac{R}{\mathfrak{p}_i}$ is a DVR for every $i$ and so by Lemma~\ref{lem:mult_sum_prim} we have that $\mult(R)=r(R)$. Therefore we can write $f=f_1 \cdots f_r$ where $f_i \in \mathfrak{n} \setminus\mathfrak{n}^2$ are irreducible that satisfy $f_i \notin \langle f_j \rangle$ for every $i \neq j$. In addition, by Remark~\ref{rem:omp} we have that $i(\mathfrak{p}_i, \mathfrak{p}_j)=1$, and so $\langle f_i, f_j \rangle = \mathfrak{n}$. Therefore, if $x$ and $y$ generate $\mathfrak{n}$, then by applying Lemma 4.2 in~\cite{greuel2019finite} with respect to $\mathfrak{n}$, for every $i$ we can find some units $a_i, b_i \in R$ such that $f_i = a_i x+b_iy$. Thus, we have that $\langle f \rangle = \langle \prod_{i=1}^{r(R)} (x+u_i y) \rangle$ for some units $u_i \in R$, and the result follows. 
\end{proof}

% If, in addition, $\chara(\kappa)$ does not divide $r(R)$, then we can write $f$ as $x^{r(R)}+y^{r(R)}$.Now, if $\chara(\kappa)$ does not divide $r(R)$, then since $R$ is strictly Henselian, there exists some $\xi \in R$ such that $\xi^{r(R)}=-1$.....

\begin{remark}
\begin{enumerate}
    \item \textup{From Proposition~\ref{prop:omp_edim2} and Theorem~\ref{thm:Morse} we can conclude that $R$ is a double point if and only if $R$ is an ordinary multiple point with $r(R)=2$. This is true since if $\mathfrak{n}=\langle x,y \rangle$ then for every units $u \neq  v \in R$ such taht $u-v \in R$ is a unit as well we have that $x+uy, x+vy$ also generate $\mathfrak{n}$.  }
    \item \textup{Proposition~\ref{prop:omp_edim2} tells us that we can intuitively think of an ordinary multiple point with $\edim(R)=2$ as "having $r(R)$ distinct tangent lines". This can be viewed as an analogue of the classical definition of an ordinary multiple point of an algebraic curve, as in Chapter 3.1 of~\cite{fulton2013intersection}.}
\end{enumerate}
\end{remark}

%If $r(R)=2$ then we are done by Theorem~\ref{thm:Morse} (noting that , where $i$ satisfies $i^2=-1$, which exists since $\chara(\kappa) \neq 2$ and $R$ is henselian, and ). Now assume that $r(R)>2$, then by the proof of Lemma~\ref{lem:ineq_edim_e} we have that $\mult\left(\frac{R}{\mathfrak{q}}\right)=r(R)-1 \geq 2$. Therefore we must have that $\mathfrak{q} \subset \mathfrak{m}^2$, and so, by Lemma~\ref{lem:ideal_edim}, we can conclude that $\edim(R) = \edim\left(\frac{R}{\mathfrak{q}}\right)$. By repeating this argument with respect to the minimal primes of $\frac{R}{\mathfrak{q}}$ (until reaching the case where $r(R)=2$), we can conclude by induction that $\edim(R)$ must be $2$. Therefore by Remark~\ref{rem:techincal_ring} we have that $\hat{R}$ is isomorphic to some quotient $\frac{S}{\langle f \rangle}$ where $(S,\mathfrak{n})$ is a two dimensional complete regular local ring and $f \in \mathfrak{n}_r$. Again, by the proof of Lemma~\ref{lem:ineq_edim_e}, we have that $R$ is of multiplicity $r(R)$ and $\frac{R}{\mathfrak{p}_i}$ is a DVR for every $i$, therefore we can write $f=f_1 \cdots f_r$ where $f_i \in \mathfrak{n}_r \setminus\mathfrak{n}^2$ are irreducible that satisfy $f_i \notin \langle f_j \rangle$ for every $i \neq j$. Therefore we can find some $x, y_1$ that generate $\mathfrak{n}_r$ and some distinct units $u_1, \dots, u_{r(R)}$ such that $f_i = x+u_iy_1$ for every $i$. .....

% \begin{remark}\label{rem:D_4}
%     \textup{From the proof of Corollary~\ref{cor:omp} we can see that if $R$ is an ordinary multiple  point with $r(R)=3$ then the completion of $R$ can be written as $\hat{R}  \cong \frac{S}{\langle xy(x+y)\rangle}$ for some generators $x$ and $y$ of some complete regular local ring $S$. Therefore, $R$ must be a $D_4$ singularity.}
% \end{remark}



%\begin{remark}
%    \textup{Note that if $(R, \mathfrak{m}, \kappa)$ is an ordinary multiple point such that $\chara(\kappa)$ does not .... (recall that $\kappa$ is algebraically closed), then we can conclude that $\hat{R}$ is isomorphic to the quotient ring $\frac{S}{\langle x^r+y^r}$ for some $2$-dimensional complete regular local ring $(S, \mathfrak{n})$ with $\mathfrak{n} = \langle x,y \rangle$ and $r(R)=r$. This is true since.......... }
%\end{remark}


\begin{definition}\label{def:A_k}
    We say that $R$ is an $A_{k}$ singularity (for $k \geq 1$) if there exists some $2-$dimensional complete regular local ring $(S, \mathfrak{n})$ such that $\mathfrak{n} = \langle x,y \rangle$ and $\hat{R} \cong \frac{S}{\langle x^2 +y^{k+1} \rangle}$. 
\end{definition}

\begin{remark}\label{rem:A_1_doublept}
    \textup{Note that if $\chara(\kappa) \neq 2$, then $R$ is a double point if and only if $R$ is an $A_1$ singularity. This is true since $\kappa$ is strictly henselian, and so we can find $i \in R$ such that $i^2=-1$. Therefore, $x^2+y^2=(x+iy)(x-iy)$ and $x,y $ generate $\mathfrak{n}$ if and only if $x+iy, x-iy$ do. }
\end{remark}


\begin{proposition}\label{cor:delta1}
    Assume that $\chara(\kappa) \neq 2$ and that $R$ is analytically unramified. Then:
    \begin{enumerate}
        \item If $\mult(R)=2$ then $R$ is an $A_k$ singularity for some $k \geq 1$.
        \item If $R$ is an $A_k$ singularity then $\mu(R)=k$
        \item $\delta(R)=1$ if and only if $R$ is either an $A_1$ or $A_2$ singularity.
        \item $\mu(R) = 2$ if and only if $R$ is an $A_2$ singularity.
    \end{enumerate} 
\end{proposition}

\begin{proof}
    For the first item, since $\mult(R)=2$ then by the first item of Remark~\ref{rem:techincal_ring} we have that $\edim(R) \leq 2$, yet since $\mult(R)=2$ then $R$ is not regular and so $\edim(R)>1$. Therefore we can conclude that $\edim(R)=2$ and so by lEMMA~\ref{lem:edim_CST}, there exists some $2-$dimensional complete regular $(S, \mathfrak{n})$ and some $f \in S$ such that $\hat{R}=\frac{S}{\langle f \rangle}$, and since $\mult(R)=2$ we can conclude that $f \in \mathfrak{n}^2 \setminus \mathfrak{n}^3$. In addition, since $R$ is analytically unramified then $\hat{R}$ is reduced and therefore so is $f \in S$. Thus, by Proposition 3.15. in~\cite{svoray2025ade}, there exists some $x,y \in S$ that generate $\mathfrak{n}$ such that $f=x^2 + y^k$ for some $k \geq 2$, and the result follows.\\

     For the second item, note that since $\mu$ is preserved under completion, it is enough to compute that $\mu\left(R\right)=k$ where $R=\frac{S}{\langle x^2 + y^{k+1}\rangle}$ for some complete regular $(S, \mathfrak{n})$ with $\mathfrak{n}= \langle x,y \rangle$. Observe that $r(R)=1$ if $k$ is even and $r(R)=2$ if $k$ is odd, since for every $n$ we have that $x^2+y^{2n+1}$ is irreducible and $x^2+ y^{2n} = (x+iy^n)(x-iy^n)$ (since $S$ is regular, therefore Henselian, and $\kappa$ is seperably closed, therefore there exists some $i\in S$ such that $i^2=-1$). If $k=2n-1$ is odd then we can directly compute that $\delta(R)=a$ using Proposition~\ref{thm:branch} (since the minimal primes of $R$ are $\langle x+iy^{n}\rangle$ and $\langle x-iy^{n}\rangle$) and if $k=2n$ is even then $\delta(R)=n$ as $\overline{R} = R\left[\frac{y}{x}, \dots, \frac{y}{x^n}\right]$.\\
    
    For the third item, if $\delta(R)=1$ then we have that $R$ is not regular, that is $\edim(R)>1$. In addition, we have that $\mu(R)=3-r(R)$. Therefore, either $\mu(R)=1$ (and by Theorem~\ref{thm:Morse} and Remark~\ref{rem:A_1_doublept} we can conclude that $R$ is an $A_1$ singularity) or $\mu(R)=2$ and $r(R)=1$.  Yet, by Lemma~\ref{cor:edim_bound} and the first item of Remark~\ref{rem:techincal_ring} we have that $\edim(R)-1 \leq \delta(R)=1$, and so we can conclude that $\edim(R)=2$. Thus we have that $\mult(R)=\edim(R)=2$ and the result follows from the previous item.\\

    For the fourth item, if $\mu(R)=2$ then by Corollary~\ref{cor:lower_bound} we have that $(r(R)-1)^2 \leq 2$, and so we must have that either $r(R)=1$ or $r(R)=2$. Yet, since $2=\mu(R)=2\delta(R)-r(R)+1$, we can conclude that $r(R)=1$. Therefore $\delta(R)=1$ and the result follows from the previous item (noting that since $\mu(R)=2$ then by Theorem~\ref{thm:Morse}  $R$ can not be an $A_1$ singularity). 
\end{proof}


\begin{remark}
\begin{enumerate}
    \item \textup{The first item of Proposition~\ref{cor:delta1} is an analogue of Theorem 2.48. in~\cite{greuel2007introduction} and the third item is analogue of Example 5.8 in~\cite{ooishi1987genera}.  }
    \item \textup{Note that Theorem 2 in~\cite{d2025upper} tells us that if $R$ is Cohen-Macaulay then $\mult(R) \leq (\edim(R)-1)\cdot\length_R\left(\frac{R}{\mathfrak{c}_R}\right)+1$ and that we have an equality if and only if $R$ is either an DVR or $\length_R\left(\frac{R}{\mathfrak{c}_R}\right)=1$. This, together with the first item of Remark~\ref{rem:conductor} and Corollary~\ref{cor:reg} gives us an alternative proof of Corollary~\ref{cor:delta1} (assuming $R$ is Cohen-Macaulay).}
\end{enumerate}
\end{remark}


% \begin{lemma}
%     Assume that $R$ is Gorenstein but not a DVR. Then $\mathfrak{c}_R \subseteq \mathfrak{m}^2$ if and only if $\delta(R) \geq \edim(R)+1$.
% \end{lemma}

% \begin{proof}
%     Short exact sequences and stuff in different cases. 
% \end{proof}

We end this section by looking at the relationship between the Milnor number of $R$ and the blow up of $R$ at the maximal ideal, known as the first neighborhood of $R$. This in fact allows us to uniquely define $\mu(R)$ based upon its numerical properties. 

\begin{definition}\label{def:blowup}
\begin{enumerate}
    \item The first neighborhood of $R$ is defined to be the (non-local) ring $R^{(1)}=\bigcup_{n \geq 0} \{a \in \overline{R} \colon a\mathfrak{m}^n \subset \mathfrak{m}^n\}$.
    \item We say that $x \in \mathfrak{m} \setminus \mathfrak{m}^2 \subset R$ is superficial if $x\mathfrak{m}^n=\mathfrak{m}^{n+1}$ for $n \gg 0$.
    \item The reduction number of $R$, denoted $\rho(R)$, is defined to be the $R-$length of the module $\frac{R^{(1)}}{R}$.
\end{enumerate} 
\end{definition}

\begin{remark}\label{rem:reduction}
\begin{enumerate}
    \item \textup{Note that superficial elements in $R$ exist since $R$ is strictly Henselian, and therefore $\kappa$ is infinite (for a proof, see for example Proposition 8.5.7 in~\cite{huneke2006integral}). In addition, if $x, \dots, x_d$ is a minimal generating sequence for $\mathfrak{m}$ in $R$ and $x \in R$ is superficial, then $R^{(1)} = R[\frac{x}{x}, \dots, \frac{x_d}{x}]$. For more details, see~\cite{northcott1956general} and Chapter 1.5 in~\cite{campillo2006algebroid}. }
    \item \textup{Since $R$ is strictly Henselian we have that for every birational finite map $R \to S$, the number of maximal ideals in $S$ is bounded above by $r(R)$ (for more information, see Lemma 10.153.3. in~\cite{stacks-project}). In particular, if $r(R)=1$ then $R^{(1)}$ is local. }
    \item \textup{The reduction number of $R$ has many relations with invariants such as the multiplicity. By~\cite{northcott1959reduction} we have that if $N \gg 0$ then $\length_R\left(\frac{R}{\mathfrak{m}^N}\right) = n \cdot \mult(R) + \rho(R)$. In addition, in~\cite{kirby1975reduction} it was shown that $\mult(R)-1 \leq \rho(R) \leq \frac{\mult(R) (\mult(R)-1)}{2}$. For more information on bounds on the reduction number of $R$, see~\cite{elias1990characterization}. }
\end{enumerate}
\end{remark}

The following Lemma is an analogue of Lemma 1.5.9 in~\cite{campillo2006algebroid}:

\begin{lemma}\label{lem:reg_blowup}
    $R^{(1)}=R$ if and only if $R$ is a DVR. 
\end{lemma}
\begin{proof}
    If $R$ is a DVR then $R$ must be normal, and so $\overline{R}=R$. But since $R \subset R^{(1)} \subset \overline{R}$ we can conclude that $R=R^{(1)}$. Now, assume that $R=R^{(1)}$ and assume towards contradiction that $R$ is not a DVR. Let $x$ be a superficial element of $R$. Since $R$ is not a DVR then $\edim(R)>1$ and so $x$ does not generate $R$. Therefore, if $y \in \mathfrak{m} \setminus \langle x \rangle$ then $\frac{y}{x} \in R^{(1)}$ but $\frac{y}{x} \notin R$, which is a contradiction. 
\end{proof}


\begin{lemma}\label{lem:rho_delta}
    If $r(R)=1$ then:
    \begin{enumerate}
        \item $\delta(R)=\rho(R)+\delta(R^{(1)})$.
        \item $\mu(R)-\mu(R^{(1)})  = 2\rho(R)$.
    \end{enumerate}
\end{lemma}

\begin{proof}
    By Remark~\ref{rem:reduction} we have that $R^{(1)}$ is local and so $\delta(R^{(1)})$ and $\mu(R^{(1)})$ are well defined. Now, since the normalization of $R^{(1)}$ is $\overline{R}$ then we have a short exact sequence
    \begin{equation*}
        0 \to \frac{R^{(1)}}{R} \to \frac{\overline{R}}{R} \to \frac{\overline{R}}{R^{(1)}} \to 0,
    \end{equation*}
    \noindent and so the result follows by applying to it length over $R$. 
\end{proof}

Since by Remark~\ref{rem:reduction}, if $r(R)=1$ then $R^{(1)}$ is local, then we can define higher neighborhoods of $R$ by inductively applying the first neighborhood of $R$. For more information on higher neighborhoods, see~\cite{northcott1955neighbourhoods, northcott1957notion}.  

\begin{definition}
    Assume that $r(R)=1$. Then the $i-$th neighborhood of $R$ is defined to be $R^{(i)}=\left(R^{(i-1)}\right)^{(1)}$, where $R = R^{(0)}$. 
\end{definition}

The following proposition is an analogue of Proposition 3.26 in Chapter I of~\cite{greuel2007introduction} and of Theorem 1.5.10 in~\cite{campillo2006algebroid}:

\begin{proposition}\label{rem:higher_NBHD}
    Assume that $r(R)=1$ and that $R$ is not regular. Then we have a sequence of inclusions $R \subsetneq R^{(1)} \subsetneq \cdots \subsetneq R^{(N)} = \overline{R}$ for some $N$ with
    \begin{equation*}
        \delta(R) = \sum_{i=1}^N \length_R \left(\frac{R^{(i)}}{R^{(i-1)}}\right).
    \end{equation*}
    In particular, we have that $\delta(R) \geq N$. 
\end{proposition}

\begin{proof}
    Since $\overline{R}$ is a Noetherian $R-$module and since by Remark~\ref{rem:reduction} we have that each $R^{(i)}$ is an $R-$submodule of $\overline{R}$, we must have that the chain of inclusions $R \subset R^{(1)} \subset \cdots $ is stationary. Therefore, there exists some $N$ such that $R^{(N)} = R^{(N+1)}$ and so Lemma~\ref{lem:reg_blowup} we have that  $R^{(N)}$ is normal, and so $R^{(N)}=\overline{R^{(N)}}=\overline{R}$. Therefore the result follows from repeatedly applying Lemma~\ref{lem:rho_delta} to the higher neighborhoods of $R$. 
\end{proof}

We end this section with a proposition that tells us that we can uniquely define the Milnor number based upon the properties we proved before. It is an analogue of Theorem 2.6 in~\cite{ploski1995milnor}. 

\begin{proposition}\label{prop:uniqueness}
    $\mu(R)$ is the unique invariant defined on the class of one-dimensional Noetherian strictly Henselian reduced Nagata local rings such 
    \begin{enumerate}
        \item $\mu(R)=0$ if and only if $R$ is a DVR, 
        \item If $r(R)=1$ then  we have $\mu(R)-\mu(R^{(1)})  = \rho(R)$, 
        \item If $\mathfrak{p}_1, \dots, \mathfrak{p}_{r(R)}$ are the minimal primes of $R$ then we have $\mu(R) -1 = \sum_{i=1}^{r(R)} \left(\mu\left(\frac{R}{\mathfrak{p}_i}\right)-1\right) + 2\sum_{i <  j} i(\mathfrak{p}_i, \mathfrak{p}_j)$.
    \end{enumerate}
\end{proposition}

\begin{proof}
    Note that $\mu(R)$ indeed satisfies these properties by Corollary~\ref{cor:lower_bound},  Corollary~\ref{cor:reg}, and Lemma~\ref{lem:rho_delta}.
    Now, assume that there is another invariant $\tilde{\mu}(R)$ that also satisfies this property. By the first property, $\tilde{\mu}(R)$ and $\mu(R)$ must agree on DVR's. By the third property, we can assume that $r(R)=1$ as $\tilde{\mu}(R)$ is uniquely defined by $\tilde{\mu}\left(\frac{R}{\mathfrak{p}_i}\right)$ and by $R$ itself.
    Therefore, by Proposition~\ref{rem:higher_NBHD} we have a sequence of inclusions of local rings $R \subset R^{(1)} \subset \cdots \subset R^{(N)} = \overline{R}$ and so $\tilde{\mu}(R)$ must equal to $\mu(R)$ by repeated application of the second property to each higher neighborhood of $R$. 
\end{proof}





% \begin{lemma}
%     If $\chara(\kappa) \neq 2$ with $r(R)=2$ and $\mult(R)=3$ then $\edim(R)=3$. 
% \end{lemma}

% \begin{proof}
%     If the minimal primes of $R$ are $\mathfrak{p}_1$ and $\mathfrak{p}_2$ then by Lemma~\ref{lem:mult_sum_prim} we can conclude that $3=\mult(R)=\mult\left(\frac{R}{\mathfrak{p}_1}\right)+ \mult\left(\frac{R}{\mathfrak{p}_2}\right)$. Therefore, we can conclude that (up to renaming the minimal primes) $\mult\left(\frac{R}{\mathfrak{p}_1}\right)=2$, and so we must have that $\mathfrak{p}_1 \subset\mathfrak{m}^2$. Therefore, by Lemma~\ref{lem:ideal_edim} we can conclude that $\edim(R)=\edim\left(\frac{R}{\mathfrak{p}_1}\right)$. Yet, since $\mult\left(\frac{R}{\mathfrak{p}_1}\right)=2$ then by the first item of Proposition~\ref{cor:delta1} we can conclude that $\frac{R}{\mathfrak{p}_1}$ is an $A_k$ singularity, which gives us that $\edim\left(\frac{R}{\mathfrak{p}_1}\right)=2$, and the result follows. 
% \end{proof}

% \begin{proposition}
%     Assume that $R$ is analytically unramified and that $\kappa$ is algebraically closed of characteristic $\neq 2$. If $r(R)=1$ and $\edim(R)=\mult(R)=3$ then the completion of $R$ is isomorphic to the quotient ring $\frac{S_3}{\langle f,g \rangle}$ where $(S_3, \mathfrak{n})$ is a $3-$dimensional complete regular local ring and $\mathfrak{n}=\langle x,y,z \rangle$ such that one of the following is true: 
% \end{proposition}







 
\section{The Semigroups $\Gamma(R)$ and $\nu(R)$}\label{sec:semigroup}

In this section we define two semigroup invariants of $R$, and we  show how they relate to $\mu(R)$. Again, we assume that $(R, \mathfrak{m}, \kappa)$ is a one-dimensional Noetherian strictly Henselian reduced Nagata local ring.

\begin{definition}
    For every $g \in \NZD(R)$ we define $l(g, R)=\length_R\left(\frac{R}{\langle g \rangle}\right) \in \mathbb{N}$. The set of all values $l(g,R)$ is denoted by $\Gamma(R)$. 
\end{definition}

\begin{remark}\label{rem:intersection}
\begin{enumerate}
    \item \textup{The motivation for $\Gamma(R)$ comes from the intersection number over algebraic curves. Specifically, given some field $k$, if $f \in k[[x,y]]$ is irreducible, then $\Gamma\left(\frac{k[[x,y]]}{\langle f\rangle}\right)$ equals to the collection of all intersection numbers $i(\langle f\rangle, \langle g\rangle)$ where $g \notin \langle f \rangle$, which is exactly the semigroup of values of $f$. For more information, see~\cite{barroso2012approach,hefez2003irreducible,castellanos2005semigroup} }
    \item \textup{Note that $\Gamma(R)$ is non trivial since for every unit $u \in R$ we have that $l(u,R)=0 \in \Gamma(R)$. In addition, by Lemma 10.63.9. in~\cite{stacks-project},  $\NZD(R)$ is exactly the complement of the union $\mathfrak{p}_1 \cup \cdots \cup \mathfrak{p}_{r(R)}$, and so is non trivial. }
\end{enumerate}
\end{remark}

The next proposition tells us that $\Gamma(R)$ is a semi-group, in addition to some properties that we use throughout this section.  

\begin{lemma}\label{lem:semigroup}
Let $g \in \NZD(R)$. Then: 
    \begin{enumerate}
        \item If $g$ is superficial then $l(g,R)=\mult(R)$. 
        \item There exists a Zariski open set $U \subset \frac{\mathfrak{m}}{\mathfrak{m}^2}$ such that $l(g,R)=\mult(R)$ for every $g+\mathfrak{m}^2 \in U$. 
        \item If  $r(R)=1$ then $l(g,R)$ equals to $\ord(g, \overline{R})= \max\{N \colon g \in \langle \pi \rangle^N\}$, where $\pi$ is a uniformizer of $\overline{R}$.
        \item If $g \in \mathfrak{m}^n$ then $l(g,R) \geq \mult(R) \cdot n + \rho(R)$.
        \item $\min(\Gamma(R)\setminus \{0\})=\mult(R)$.
        \item Given $g, h \in \NZD(R)$ we have that $l(gh, R)=l(g, R)+l(h,R)$.  
        %\item If $n > m(R)$ then $n \in \Gamma(R)$. 
    \end{enumerate}
\end{lemma}

\begin{proof}
The first and second item follow from Theorem 8.6.6 and Proposition 11.2.2 in~\cite{huneke2006integral}, respectively (with the first item also being based upon~\cite{northcott1954reductions, trung2003constructive}). In addition, the third item follows from Example A.3.1 in~\cite{fulton2013intersection}. The fourth item follows from Remark~\ref{rem:reduction}, noting that if $g \in \mathfrak{m}^r$ then $l(g,R) \geq \length_R\left(\frac{R}{\mathfrak{m}^r}\right)$.  \\

For the fifth item, it is enough to show that $\length_R\left(\frac{R}{\langle g^n \rangle }\right) \geq n\cdot l(g,R)$, since as $g^n \in \mathfrak{m}^n$ for every $n$, we can conclude that 
\begin{equation*}
    l(g,R) = \frac{n \cdot \length_R\left(\frac{R}{g}\right)}{n} \geq \frac{\length_R\left(\frac{R}{g^n}\right)}{n} \geq   \frac{\length_R\left(\frac{R}{\mathfrak{m}^n}\right)}{n} \to \mult(R). 
\end{equation*}

\noindent We show this by induction, noting that the case $n=1$ is true from the definition of $l(g,R)$. Now, assuming it is true for $n-1$. Since $g \in (\langle g^n \rangle \colon \langle g^{n-1}\rangle)$, then by looking at the short exact sequence 

\begin{equation*}
    0 \to \frac{R}{(\langle g^n \rangle \colon \langle g^{n-1}\rangle)} \to \frac{R}{\langle g^n\rangle} \to \frac{R}{\langle g^{n-1}\rangle} \to 0,
\end{equation*}

\noindent we have that $\length_R\left(\frac{R}{\langle g^n\rangle}\right) = \length_R\left(\frac{R}{\langle g^{n-1}\rangle}\right) +\length_R\left(\frac{R}{(\langle g^n \rangle \colon \langle g^{n-1}\rangle)}\right)  \geq (n-1) \cdot l(g,R) + l(g,R) = n\cdot l(g,R)$, and the result follows. \\

For the sixth item, it is enough to show that $\frac{R}{\langle h \rangle} \cong \frac{\langle g \rangle}{\langle gh \rangle}$, as this isomorphism gives us the short exact sequence 

\begin{equation*}
    0 \to \frac{R}{\langle h \rangle} \to \frac{R}{\langle gh \rangle} \to \frac{R}{\langle g \rangle} \to 0.
\end{equation*}

\noindent In order to prove the isomorphism, we look at the map $R \to \frac{\langle g \rangle}{\langle gh \rangle}$ defined by $f \mapsto fg$. Assume that $fg \in \langle gh \rangle$. Then there exists some $a \in R$ such that $fg = gha$, and so $g(f-ha)=0$. Since $g \in \NZD(R)$, we must have that $f=ha$, and so $f \in \langle h \rangle$. This gives us that the kernel of the map is exactly $\langle h \rangle$, and the result follows. 
\end{proof}

\begin{remark}\label{rem:semigroup}
\begin{enumerate}
    \item \textup{In Section 2 of~\cite{svoray2025detecting}, an analogue of the Milnor and Tjurina numbers is defined over a mixed characteristic DVR as the multiplicity of a specific one dimensional ring. Their definitions, which is based upon Proposition 3.3. in that paper, uses an analogue of items 1 and 2 of Lemma~\ref{lem:semigroup}, as item 2 tells us that $l(g,R)$ equals to the multiplicity of $R$ for a "generic element in $\mathfrak{m} \setminus \mathfrak{m}^2$".}
    \item \textup{Items 4 and 5 in Lemma~\ref{lem:semigroup} give us analogous results to those of the intersection number of algebraic curves (as in Remark~\ref{rem:intersection}). For more information, see Section 3.1 of Chapter I in~\cite{greuel2007introduction}.  }
    \item     \textup{The third item in Lemma~\ref{lem:semigroup} can be viewed as an analogue of Proposition 1.1 in~\cite{ploski1995milnor} and of Proposition 3.2 in Chapter I of ~\cite{greuel2007introduction}. Specifically, if $r_R=1$ we can think of the normalization map $R \to \overline{R}$ as a generalization of the idea of "parametrization of a curve", as it allows us to write the maximal ideal of $R$ as generated by powers of the uniformizer of $R$. Fore more on parametrization of curves, see Section 3.4 of Chapter I in~\cite{greuel2007introduction} and Chapter 1 of~\cite{campillo2006algebroid}.}
    \item \textup{As mentioned above, item 6 of Lemma~\ref{lem:semigroup} tells us that $\Gamma(R) \subset \mathbb{N}$ is a semigroup, i.e. contains zero and closed under addition. In fact, as we in Proposition~\ref{prop:semigrp}, $\Gamma(R)$ is a numerical semigroup, i.e $\gcd(\Gamma(R))=1$, which is in fact equivalent to having that $N \in \Gamma(R)$ for every $N \gg 0$ (for a proof, see Proposition 1.1 in~\cite{zariski2006moduli}). Given a numerical semigroup $S \subset \mathbb{N}$, the smallest $c$ which has the property that  if $N \geq c$ then $N \in S$ is called the conductor of $S$, and as we see in Proposition~\ref{prop:mu_vec} and Corollary~\ref{cor:Gamma_cond}, the conductor of $R$ (as defined in the first item of Remark~\ref{rem:conductor}), plays an analogous role for the conductor of $\Gamma(R)$. In fact, the conductor of any semigroup must be an even number. For more information on semigroups and their relation to curves and rings, see~\cite{assi2020numerical} or Section 1.1 in~\cite{barroso2012approach}.}
\end{enumerate}
\end{remark}


% The following Proposition is a general result on sub-semigroups of $\mathbb{N}$ (i.e. sets $S \subset \mathbb{N}$ that contain $0$ and are closed under addition) that we use throughout this section. For more details on such semigroups, see XXX. 

% \begin{proposition}\label{prop:semigroup_prop}
%     Let $S \subset \mathbb{N}$ be a semigroup. Then:
%     \begin{enumerate}
%         \item $S$ is finitely generated (i.e. there exists some $x_1, \dots, x_d$ such that $S = \{n_1x_1 + \cdots + x_dx_d \colon n_i \in \mathbb{N}\}$. 
%         \item There exists a minimal constant $c(S)$ such that for every for every $N \geq c(S)$ that is divisible by $\gcd(S)$ we have that $N \in S$.
%         \item 
%     \end{enumerate}
% \end{proposition}

% \begin{proof}
%     For the first item, let $Y \subset S$ be a finite set such that $\gcd(Y)=\gcd(S)$. Denote by $X$ the smallest semigroup in $\mathbb{N}$ containing $Y$. Then since $Y \subset S$ we have that $X \subset S$. Since $Y$ is a finite set, then we can find some $x\in X$ such that if $n \geq x$ is in $X$ then $n=x+k\gcd(S)$ for some $k$. Therefore, all numbers that are larger than $x$ and divisible by $\gcd(S)$ are in $X$. This in particular gives us that all elements in $S$ that are greater than $x$ are in $X$. Therefore we can conclude that every element in $S$ can be written as a sum of elements in $X$ and elements in $S   \cap \{n \in \mathbb{N} \colon n \leq x \}$. \\

%     For the second item, 
% \end{proof}

\begin{corollary}\label{Cor:DVR_Gamma}
    $R$ is a DVR if and only if $\Gamma(R)=\mathbb{N}$. 
\end{corollary}
\begin{proof}
    By Lemma~\ref{lem:semigroup} we have that $\Gamma(R)=\mathbb{N}$ if and only if $\mult(R)=1$, which is equivalent to $R$ being regular. 
\end{proof}

\begin{remark}
    \textup{Note that if $r(R)=1$ then by Proposition~\ref{rem:higher_NBHD} we have a sequence of inclusions $R \subset R^{(1)} \subsetneq \cdots \subsetneq R^{(N)} = \overline{R}$, and so by Lemma~\ref{lem:semigroup} and Corollary~\ref{Cor:DVR_Gamma}  we have a sequence of inclusions $\Gamma(R) \subsetneq \Gamma(R^{(1)}) \subsetneq \cdots \subsetneq \Gamma(R^{(N)}) = \mathbb{N}$ (as the delta invariant decreases under blow up by Lemma~\ref{lem:rho_delta}). In particular, we can conclude that $\mult(R) \geq \mult(R^{(1)}) \geq \cdots \geq \mult(R^{(N)}) = 1$. }
\end{remark}



\begin{lemma}\label{lem:l_irred}
     Assume that $\mathfrak{p}_1, \dots, \mathfrak{p}_{r(R)}$ are the minimal primes of $R$. Then for every $g \in \NZD(R)$ we have that $l(g,R)=\sum_{i=1}^{r(R)} l(g_i, \frac{R}{\mathfrak{p}_i})$, where $g_i$ is the image of $g$ under the projection $R \to \frac{R}{\mathfrak{p}_i}$.
\end{lemma}

\begin{proof}
    Note that since $R \to \frac{R}{\mathfrak{p}_i}$ is a surjective map, then by Lemma 10.52.5. in~\cite{stacks-project} we have that $l(g_i, \frac{R}{\mathfrak{p}_i}) = \length_R \left(\frac{R}{\mathfrak{p}_i + \langle g \rangle}\right)$. As in Lemma~\ref{lem:mult_sum_prim}, denote $\mathfrak{q}=\bigcap_{i=1}^{r(R)-1} \mathfrak{p}_i$ and recall that $\mathfrak{q}\cap \mathfrak{p}_{r(R)} = \{0\}$. Now, we have a short exact sequence 
    \begin{equation*}
        0 \to R \to \frac{R}{\mathfrak{q}} \oplus \frac{R}{\mathfrak{p}_{r(R)} } \to \frac{R}{\mathfrak{q}+\mathfrak{p}_{r(R)} } \to 0,
    \end{equation*}
     \noindent and so by applying to it $-\bigotimes_R \frac{R}{\langle g \rangle}$ we can have an exact sequence 
    \begin{equation*}
         0 \to  \Tor^1_R\left( \frac{R}{\mathfrak{q}+\mathfrak{p}_{r(R)} }, \frac{R}{\langle g \rangle} \right) \to  \frac{R}{\langle g \rangle} \to \frac{R}{\mathfrak{q} + \langle g \rangle} \oplus \frac{R}{\mathfrak{p}_{r(R)} + \langle g \rangle} \to \frac{R}{\mathfrak{q}+\mathfrak{p}_{r(R)} + \langle g \rangle} \to 0. 
    \end{equation*}
    \noindent Yet, we have that $\Tor^1_R\left( \frac{R}{\mathfrak{q}+\mathfrak{p}_n}, \frac{R}{\langle g \rangle} \right)$ is exactly $( 0 \colon g) = \{ x \in \frac{R}{\mathfrak{q}+\mathfrak{p}_n} \colon gx =0\}$. Since $\frac{R}{\mathfrak{q}+\mathfrak{p}_n}$ is an Artinian local ring, so by looking at the short exact sequence
    \begin{equation*}
        0 \to (0 : g) \to \frac{R}{\mathfrak{q}+\mathfrak{p}_n} \xrightarrow{x \mapsto gx}  \frac{R}{\mathfrak{q}+\mathfrak{p}_n}  \to  \frac{R}{\mathfrak{q}+\mathfrak{p}_n+ \langle g \rangle} \to 0,
    \end{equation*}
    \noindent we can conclude that the $R-$length of $\frac{R}{\mathfrak{q}+\mathfrak{p}_n+ \langle g \rangle} $ is the same as the $R-$length $\Tor^1_R\left( \frac{R}{\mathfrak{q}+\mathfrak{p}_n}, \frac{R}{\langle g \rangle} \right)$. Therefore we have that $\length_R \left ( \frac{R}{\langle g \rangle} \right) = \length_R \left(\frac{R}{\mathfrak{q} + \langle g \rangle}\right) +  \length_R \left( \frac{R}{\mathfrak{p}_n+ \langle g \rangle}\right)$, and the result follows by repeating the argument with respect to $\mathfrak{q}$. 
\end{proof}

%As in Lemma~\ref{lem:mult_sum_prim}, denote $\mathfrak{q}=\bigcap_{i=1}^{n-1} \mathfrak{p}_i$ and since $\mathfrak{q}\cap \mathfrak{p}_n = \{0\}$ we have that $\mathfrak{p}_n\mathfrak{q}=\{0\}$. Therefore we have that $l(g,R)=\length_R \left( \frac{R}{\langle g \rangle + \mathfrak{p}_n\mathfrak{q}}\right)$. Now, similar to the the fifth item in Lemma~\ref{lem:semigroup}, we show that $\length_R \left( \frac{R}{\langle g \rangle + \mathfrak{p}_n\mathfrak{q}}\right) = \length_R \left( \frac{R}{\langle g \rangle + \mathfrak{q}}\right) + \length_R \left( \frac{R}{\langle g \rangle + \mathfrak{p}_n}\right)$, which ends the proof as we can then preform the same argument with respect to the minimal ideals of $ \frac{R}{\mathfrak{q}}$. Since we have a short exact sequence 
    % \begin{equation*}
    %     0 \to \frac{\langle g \rangle + \mathfrak{p}_n}{\langle g \rangle + \mathfrak{p}_n\mathfrak{q}} \to \frac{R}{\langle g \rangle + \mathfrak{p}_n\mathfrak{q}} \to \frac{R}{\langle g \rangle + \mathfrak{p}_n} \to 0, 
    % \end{equation*}
    % \noindent we claim that we have an isomorphism $\frac{\langle g \rangle + \mathfrak{p}_n}{\langle g \rangle + \mathfrak{p}_n\mathfrak{q}} \cong \frac{R}{\langle g \rangle + \mathfrak{q}}$. Since we have an inclusion $\langle g \rangle + \mathfrak{p}_n \subset R$, we get a projection map $\langle g \rangle + \mathfrak{p}_n \to \frac{R}{\langle g \rangle + \mathfrak{q}}$

\begin{remark}\label{cor:sum_length}
\begin{enumerate}
    \item \textup{Note from Remark~\ref{rem:semigroup} and Lemma~\ref{lem:l_irred} we can conclude that if $g \in \NZD(R)$ then $l(g,R) = \sum_{i=1}^{r(R)} \ord\left(g_i, \overline{\frac{R}{\mathfrak{p}_i}}\right)$.}
    \item \textup{We can use Lemma~\ref{lem:l_irred} to give an alternative proof to Lemma~\ref{lem:mult_sum_prim}. As in the first item of Lemma~\ref{lem:semigroup}, We can find sets $U, U_1, \dots, U_{r(R)} \subset \mathfrak{m}$ that are Zariski open sets in $\frac{\mathfrak{m}}{\mathfrak{m}^2}$ and $\frac{\mathfrak{m}}{\mathfrak{m}^2 + \mathfrak{p}_i}$ for every $i$ (respectively). Since they are Zariski open set, their intersection must be non-empty. Thus, we can find some $r \in R$ such that $l(r,R)=\mult(R)$ and $l\left(r_i,\frac{R}{\mathfrak{p}_i}\right)=\mult\left(\frac{R}{\mathfrak{p}_i}\right)$, and so the result follows.} 
\end{enumerate}
\end{remark}

%Therefore, since $\overline{R}= \prod_{i=1}^{r(R)} \overline{\frac{R}{\mathfrak{p}_i}}$, then we can find for every $i$ a uniformizer $\pi_i$ of $\overline{\frac{R}{\mathfrak{p}_i}}$ such that $g = (\pi_1^{N_1}, \dots, \pi_{r(R)}^{N_{r(R)}})$ where $N_i= \ord\left(g_i, \overline{\frac{R}{\mathfrak{p}_i}}\right)$ for every $i$. 



%Note that Lemma~\ref\label{lem:semigroup} tells us that $\Gamma(R)$ is a subsemigroup of $\mathbb{N}$, that is, contains $0$ and is closed under addition. Therefore we can conclude that $\Gamma(R)$ is finitely generated as a semigroup and agrees $\gcd(\Gamma(R)) \mathbb{N}$ for big enough values (in fact, we show that $\gcd(\Gamma(R))=1$). We formalize this idea using the definition of the conductor of $R$. For more information on semigroups, see XXX. 

% \begin{definition}
%     The conductor of $R$, denoted by $c(R)$, is the smallest integer such that for every $N \geq c(R)$ we have that 
% \end{definition}


% \begin{remark}
%     \textup{Inspired by Proposition~\ref{prop:gcd_Gamma_irred}, one would expect that $\gcd(\Gamma(R))=r(R)$. Yet this is not true since $\gcd(\Gamma(R))$ must divide $\mult(R) \in \Gamma(R)$, yet $r(R)$ need not divide $\mult(R)$ (for example, if $r(R)=2$ but $\mult(R)=3$, as in Proposition~\ref{prop:E6Dk}.}
% \end{remark}


We now turn to defining $\nu(R)$,  a version of $\Gamma(R)$ that is a subset of $\mathbb{N}^{r(R)}$, which can be thought of as a "component-wise" version of $\Gamma(R)$ based upon the minimal primes of $R$. The semigroup $\nu(R)$ has been studied in the past, both for a reduced algebraic curve over a field with multiple branches and for one dimensional local ring. For more information, Section 3 of Chapter I in~\cite{greuel2007introduction} and~\cite{de1987semigroup,campillo1994gorenstein, d2025value, zariski2006moduli, campillo1994gorenstein}. 

\begin{definition}\label{def:nu}
    $\nu(R)$ is the set of all tuples $\nu(g) = \left( l\left(g_i,\frac{R}{\mathfrak{p}_i}\right) \right)_{i=1}^{r(R)} \in \mathbb{N}^{r(R)}$ for every $g \in \NZD(R)$ (where $g_i$ is the image of $g$ under the map $R \to \frac{R}{\mathfrak{p}_i}$). 
\end{definition}

The following lemma summarizes some important properties that $\nu(R)$ has:

\begin{proposition}\label{prop:nu_beasic_prop}
    Let $\mathfrak{p}_1, \dots, \mathfrak{p}_{r(R)}$ be the minimal primes of $R$. Then we have that:
    \begin{enumerate}
        \item $(N_1, \dots, N_{r(R)}) \in \nu(R)$ if and only if for every $i$ we can find a uniformizer $\pi_i$ of $\overline{\frac{R}{\mathfrak{p}_i}}$ and a unit $u_i \in \overline{\frac{R}{\mathfrak{p}_i}}$ such that $(u_1\pi_1^{N_1}, \dots, u_{r(R)}\pi_{r(R)}^{N_{r(R)}}) \in R \subset \overline{R} = \prod_{i=1}^{r(R)} \overline{\frac{R}{\mathfrak{p}_i}}$.
        \item The minimal tuple in $\nu(R) \setminus \{(0, \dots, 0)\}$ is $\left(\mult\left(\frac{R}{\mathfrak{p}_1}\right), \dots, \mult\left(\frac{R}{\mathfrak{p}_{r(R)}}\right)\right)$. 
        \item $\Gamma(R)$ is the image of $\nu(R)$ under the linear transformation $s \colon \mathbb{N}^n \to \mathbb{N}$ defined by $s(x_1,\dots, x_n)=x_1+\cdots+x_n$.
        \item The image of $\nu(R)$ under the projection $\mathbb{N}^{r(R)} \to \mathbb{N}^{r(R)-1}$ defined by $(N_1, \dots, N_{r(R)}) \mapsto (N_1, \dots, N_{i-1}, N_{i+1}, \dots,  N_{r(R)})$ is exactly $\nu\left(\frac{R}{\mathfrak{q}} \right)$ where $\mathfrak{q}=\bigcap_{j \neq i} \mathfrak{p}_j$. 
        \item The image of $\nu(R)$ under the projection on the $i-$th coordinate is exactly $\Gamma\left(\frac{R}{\mathfrak{p}_i}\right)$. 
    \end{enumerate}
\end{proposition}

\begin{proof}
    These follow directly from Remark~\ref{rem:techincal_ring},  Lemma~\ref{lem:semigroup}, Lemma~\ref{lem:l_irred}, and Remark~\ref{cor:sum_length}. 
\end{proof}


\begin{remark}\label{rem:Gam_comp}
\begin{enumerate}
    \item \textup{Note that $\Gamma(R)=\Gamma(\hat{R})$ and $\nu(R) =\nu(\hat{R})$, since if $r(R)=1$ this follows from Remark~\ref{rem:semigroup} as the order of an element does not change under completion, and the general case follows from Proposition~\ref{prop:nu_beasic_prop}.   }
    \item \textup{From Lemma~\ref{lem:l_irred} and Proposition~\ref{prop:nu_beasic_prop} we can conclude that given $N_i \in \Gamma\left(\frac{R}{\mathfrak{p}_i}\right)$ for some $i$ then we can find some $N_1, \dots, N_{i-1}, N_{i+1}, \dots, N_{r(R)}$ and a uniformizer $\pi_i$ of $\overline{\frac{R}{\mathfrak{p}_i}}$ for every $i$ such that $(\pi_1^{N_1}, \dots, \pi_{r(R)}^{N_{r(R)}}) \in R$. Therefore, we have that $(N_1, \dots, N_{r(R)}) \in \nu(R)$ and so $N_1 + \cdots + N_{r(R)} \in \Gamma(R)$.}
    \item \textup{Note that the intersection of $\nu(R)$ with any of the axes must be the origin. This true since if $(N_1, \dots, N_{r(R)})\in \nu(R)$ such that $N_i=0$ for some $i$ then by Proposition~\ref{prop:nu_beasic_prop}, we some $g$ such that $l(g_i, \frac{R}{\mathfrak{p}_i})=0$, and so $g_i$ is a unit. But since $g_i$ is the image of $g$ under the map $R \to \frac{R}{\mathfrak{p}_i}$, we can conclude that $g$ must be a unit and so $\nu(g)=(0,\dots,0)$. }
\end{enumerate}
\end{remark}


\begin{lemma}\label{lem:ord_ideal}
    Given $g,h \in R$, if $\nu(g)=\nu(h)$ then $f=gu$ for some unit $u \in \overline{R}$. 
\end{lemma}

\begin{proof}
    Since $\overline{R}= \prod_{i=1}^{r(R)} \overline{\frac{R}{\mathfrak{p}_i}}$ then we can write $g=(g_1, \dots, g_{r(R)})$ and $h=(h_1, \dots, h_{r(R)})$ as elements in $\overline{R}$. By Lemma~\ref{lem:semigroup} we have that $N_i= \ord(g_i, \overline{\frac{R}{\mathfrak{p}_i}}) = \ord(h_i, \overline{\frac{R}{\mathfrak{p}_i}})$. Therefore, if $\pi_i$ is a uniformizer of $\overline{\frac{R}{\mathfrak{p}_i}}$ then we can conclude that $g_i = u_i \pi_i^{N_i}$ and $h_i = v_i \pi_i^{N_i}$ for some units $u_i, v_i \in \overline{\frac{R}{\mathfrak{p}_i}}$, and so $\langle g_i \rangle = \langle h_i \rangle$. Therefore, as elements in $\overline{R}$, we have that $\langle g \rangle = \langle (\pi_1^{N_1}, \dots, \pi_{r(R)}^{N_{r(R)}})\rangle = \langle h \rangle$, and the result follows. 
\end{proof}

\begin{remark}
    \textup{In Lemma~\ref{lem:ord_ideal}, we can not conclude that $u\in R$. That is, if two elements in $R$ generate the same ideal in $\overline{R}$ we can not conclude that they generate the same ideal in $R$. For example, let $R=\frac{k[[x,y]]}{\langle x^2 - y^3 \rangle}$ for some field $k$, then $\langle x \rangle \neq \langle x+y\rangle$ but the normalization of $R$ is $R \cong k[[t^2, t^3]] \subset k[[t]] = \overline{R}$, in which we have that $y \mapsto t^2$ and $x+y \mapsto t^2(1+t)$. }
\end{remark}


We use the definition of $\nu(R)$ together with Proposition~\ref{prop:nu_beasic_prop} to show that $\Gamma(R)$ is a numerical semigroup:

\begin{proposition}\label{prop:semigrp}
    $\gcd(\Gamma(R))=1$.
\end{proposition}

\begin{proof}
    Recall from the fourth item of Remark~\ref{rem:semigroup} that $\gcd(\Gamma(R))=1$ if and only if $N\in \Gamma(R)$ for $N\gg0$. We first prove this in the case where $r(R)=1$. By Remark~\ref{rem:semigroup} we have that $l(r,R)$ is the order of $r$ in $\overline{R}$. let $\pi$ be a uniformizer of $\overline{R}$, which is a DVR since $r(R)=1$. Since $\Frac(R)=\Frac(\overline{R})$, we can find some $g,h \in R$ such that $\frac{g}{h}=\pi \in \Frac(R)$. Therefore we can conclude that $\ord(g, \overline{R})-\ord(h,\overline{R})=1$ and so $\gcd(l(g,R), l(h,R))=1$. Now, if $r(R)>1$ denote the minimal primes of $R$ by $\mathfrak{p}_1, \dots, \mathfrak{p}_{r(R)}$. By the $r(R)=1$ case
    we have that for every $i$ there exists some $M_i$ such that $N \in \Gamma\left(\frac{R}{\mathfrak{p}_i}\right)$ for every $N > M_i$. Therefore, given a tuple $(N_1, \dots, N_{r(R)}) \in \mathbb{N}^{r(R)}$ such that $N_i > M$ for every $i$, we can find for every $i$ some $g_i \in \frac{R}{\mathfrak{p}_i}$ such that $\ord\left(g_i, \overline{\frac{R}{\mathfrak{p}_i}}\right)=N_i$. Therefore, from the second item of Remark~\ref{rem:Gam_comp} we have that $g=(g_1, \dots, g_{r(R)}) \in R$ satisfies $\nu(g) = (N_1, \dots, N_{r(R)})$. So, for every $N \gg 0$, we can find some $(N_1, \dots, N_{r(R)})$ such that $N_1 + \cdots N_{r(R)}=N$ and $N_i > M$ for every $i$, which would give us that $N \in \Gamma(R)$ by Proposition~\ref{prop:nu_beasic_prop}. 
\end{proof}

As mentioned in Remark~\ref{rem:semigroup}, since $\Gamma(R)$ is a semigroup, it has a conductor $c$. In  order to compute this element $c$, we first have to define the relative Milnor number of $R$, inspired by Section 3 of~\cite{ploski1995milnor}. 

\begin{definition}\label{def:RelMIlnor}
    If $\mathfrak{p}_1, \dots, \mathfrak{p}_r$ are the minimal primes of $R$, we define the $i-$th relative Milnor number of $R$ to be $\mu_i(R)=\mu\left(\frac{R}{\mathfrak{p}_i}\right) + \sum_{j \neq i} i(\mathfrak{p}_i, \mathfrak{p}_j)$. In addition, we define the vector of relative Milnor numbers of $R$ to be $\vec{\mu}(R)=(\mu_1(R), \dots, \mu_r(R))$. 
\end{definition}

\begin{lemma}\label{lem:vec_mu_sum}
    If $\mathfrak{p}_1, \dots, \mathfrak{p}_r$ are the minimal primes of $R$ then $\mu_1(R)+\cdots+\mu_r(R)=\mu(R)-r(R)+1$. 
\end{lemma}

\begin{proof}
    This follows directly from Corollary~\ref{cor:lower_bound}. 
\end{proof}


\begin{proposition}\label{prop:mu_vec}
    Assume $R$ is Gorenstein. Then $\vec{\mu}(R)$ is the minimal element in $\nu(R) \subset \mathbb{N}^{r(R)}$ such that $\vec{\mu}(R) + \mathbb{N}^{r(R)} \subset \nu(R)$. 
\end{proposition}

\begin{proof}
 Since $\overline{R} = \prod_{i=1}^{r(R)} \overline{\frac{R}{\mathfrak{p}_i}}$ is a finite product of DVRs, it must be a PIR (i.e. every non-trivial ideal is principal), and so if $\pi_i$ is a uniformizer of $\overline{\frac{R}{\mathfrak{p}_i}}$ for every $i$, then by Remark~\ref{rem:semigroup} we can conclude that $\mathfrak{c}_R$, viewed as an ideal of $\overline{R}$ is generated by $(\pi_1^{\mu_1(R)}, \dots, \pi_{r(R)}^{\mu_{r(R)}(R)})$. Therefore we have that $\vec{\mu}(R) \in \nu(R)$ since $(\pi_1^{\mu_1(R)}, \dots, \pi_{r(R)}^{\mu_{r(R)}(R)}) \in R$. Since $\mathfrak{c}_R$ is the conductor of $\overline{R}$ over $R$, we have that $(\pi_1^{\mu_1(R)}, \dots, \pi_{r(R)}^{\mu_{r(R)}(R)}) \cdot \overline{R} \subset R$. Thus, from Proposition~\ref{prop:nu_beasic_prop}, we can conclude that for every $(\alpha_1, \dots, \alpha_{r(R)}) \leq \vec{\mu}(R)$ we must have that $(\pi_1^{\alpha_1}, \dots, \pi_{r(R)}^{\alpha_{r(R)}}) \cdot \overline{R} \not\subset R$. Therefore we have that $\vec{\mu}(R) \in \nu(R)$ and given some $\vec{\beta}=(\beta_1, \dots, \beta_{r(R)}) \in \mathbb{N}^{r(R)}$ we have that $(\pi_1^{\mu_1(R)+\beta_1}, \dots, \pi_{r(R)}^{\mu_{r(R)}(R)+\beta_{r(R)}}) \in R$ and so $\vec{\beta}+\vec{\mu}(R) \in \nu(R)$, while no tuple $\vec{\alpha} < \vec{\mu}(R)$ satisfies this property. 
\end{proof}

\begin{corollary}\label{cor:Gamma_cond}
    Assume $R$ is Gorenstein. Then the conductor of $\Gamma(R)$ is smaller than $2\delta(R) = \mu(R)-r_R+1$ with equality if $r(R)=1$. 
\end{corollary}

\begin{proof}
    This follows directly from Proposition~\ref{prop:mu_vec}.
\end{proof}


\begin{remark}\label{rem:mu-r(R)-gamma}
\begin{enumerate}
    \item \textup{From Corollary~\ref{cor:Gamma_cond} we can in fact conclude that if $r(R)=1$ then $\mu(R)-1=2 \delta(R)-1 \notin \Gamma(R)$, as otherwise it would be the minimal element. In addition, from Proposition 1.5 in~\cite{barroso2012approach} we can conclude that the number of integers in $\mathbb{N} \setminus \Gamma(R)$ is exactly $\delta(R)$.}
    \item \textup{Note that the conductor of $\Gamma(R)$ need not be $2\delta(R)$ if $r(R)>1$, as it is possible that there exists some tuple $(N_1, \dots, N_{r(R)}) \in \Gamma(R)$ with $N_1 + \cdots + N_{r(R)} = 2\delta(R)-1$  but $(N_1, \dots, N_{r(R)})$ would not satisfy the condition in Proposition~\ref{prop:mu_vec}. For a concrete example, see the computation for $E_7$ in Lemma~\ref{lem:ADE_nu_grp}.}
    \item \textup{From Theorem 4.8 in~\cite{campillo1994gorenstein} (which in turn is based upon~\cite{de1988gorenstein}) we can conclude another  property of $\vec{\mu}(R)$: Given some $ (\alpha_1, \dots, \alpha_r) \in \mathbb{N}^r$, we set $\Delta_i(\alpha_1, \dots, \alpha_r)$ to be the set of all tuples $(\beta_1, \dots, \beta_r) \in \mathbb{N}^r$ such that $\beta_i = \alpha_i$ but $\beta_j >\alpha_j$ for every $j \neq i$. Then $R$ is Gorenstein with minimal primes $\mathfrak{p}_1, \dots, \mathfrak{p}_{r(R)}$ is equivalent to $\nu(R)$ satisfying that $(\alpha_1, \dots, \alpha_{r(R)}) \in \nu(R)$ if and only if $\Delta_i\left(\mu_1(R)-1-\alpha_1, \dots, \mu_{r(R)}(R)-1-\alpha_{r(R)} \right) \cap \nu(R) = \emptyset$ for every $i$. }
\end{enumerate}
%\begin{enumerate}

%    \item \textup{We can provide a proof of Proposition~\ref{prop:mu_vec} in the case $r(R)=1$ based upon Remark~\ref{rem:semigroup}. Since $r(R)=1$ then $\mu(R)$ is in fact the $R-$length of $\frac{\overline{R}}{\mathfrak{c}_R}$ (as in Remark~\ref{rem:conductor}). Therefore, if $\pi$ is a uniformizer of $\overline{R}$, then by viewing $R \subset \overline{R}$, then $\mu(R)$ satisfies $\pi^{\mu(R)} \cdot \overline{R} \subset R$ and $\pi^{\mu(R)} \cdot \overline{R} \not\subset R$. Therefore we have that $\pi^{\mu(R)} \in R$ but $\pi^{\mu(R)-1} \notin R$. Thus, from Remark~\ref{rem:semigroup} we can conclude that $\mu(R) \in \Gamma(R)$ but $\mu(R)-1\notin \Gamma(R)$. Now, if $N > \mu(R)$ then since $\pi^{\mu(R)} \cdot \overline{R} \subset R$ we have that $\pi^{N} = \pi^{\mu(R)} \cdot \pi^{N-\mu(R)} \in R$, and so $N \in \Gamma(R)$. }
%\end{enumerate}
\end{remark}

We can use Proposition~\ref{prop:mu_vec} to give a version of Theorem~\ref{thm:Morse} based upon the semigroup $\Gamma(R)$. 

\begin{corollary}\label{cor:A_1_Gamma}
    Assume $R$ is a Gorenstein and analytically unramified that is not a domain. Then $\Gamma(R) = \{0\} \cup (2 + \mathbb{N})$ if and only if $R$ is a double point.
\end{corollary}

\begin{proof}
    First, if $R$ is a double point then $\hat{R} = \frac{S}{\langle xy \rangle}$ where $(S, \mathfrak{n})$ is some complete regular $2-$dimensional ring with $\mathfrak{n}=\langle x,y \rangle$. Let $g \in \NZD(R)$ be a non unit, then there exists some $m,n > 1$ and some units $u,v \in S$ such that $g=ux^n+vy^m + \langle xy \rangle$. Therefore we have that the $R-$length of $\frac{S}{\langle xy, ux^n+vy^m \rangle}$ is $n+m$, and so by Remark~\ref{rem:Gam_comp} we can conclude that $\Gamma(R) = \{0, 2, 3, \dots \}$. \\

    Second, if $\Gamma(R) = \{0, 2,3, \dots\}$ then from Proposition~\ref{prop:mu_vec} together with the fourth item of Lemma~\ref{lem:semigroup} we can conclude that $2\delta(R)=\mult(R)$. Yet, from Lemma~\ref{cor:edim_bound} we have that $2\delta(R) \geq 2\mult(R)-2$ and so we must have that $2 \geq \mult(R)$. Since $\mult(R)=1$ if and only if $\Gamma(R)=\mathbb{N}$, we can conclude that $\mult(R)=2$ and so $\delta(R)=1$. Since $R$ is not a domain we have that $r(R)>1$ and so by Corollary~\ref{cor:mult_r} we must have that $r(R)=2$. Therefore we can conclude that $\mu(R)=1$ and the result follows Theorem~\ref{thm:Morse}.
\end{proof}

\begin{remark}\label{rem:A_12groups}
    \textup{If we remove the assumption in Corolalry~\ref{cor:A_1_Gamma} that $R$ is not a domain, then we can conclude that $\Gamma(R)= \{0\} \cup (2 + \mathbb{N})$ if and only if $\delta(R)=1$, which by Proposition~\ref{cor:delta1} is equivalent (if we assume in addition that $\chara(\kappa) \neq 2$) to $R$ being either an $A_1$ or an $A_2$ singularity. This shows us that the semigroup $\Gamma(R)$ does not detect the number of minimal primes that $R$ has. }
\end{remark}

We end this section by showing that we can also use the semigroup $\Gamma(R)$ to prove a version of Noether's theorem for $\Gamma(R)$, as an analogue of Theorem 3.1 and Theorem 3.2 in~\cite{ploski1995milnor}:

\begin{proposition}
    Let $g,h \in \NZD(R)$ and let $\mathfrak{p}_1, \dots, \mathfrak{p}_{r(R)}$ be the minimal primes of $R$. Then:
    \begin{enumerate}
        \item If $l(h, \frac{R}{\mathfrak{p}_i}) \geq l(g, \frac{R}{\mathfrak{p}_i}) + \mu_i(R)$ for every $i$ then $h \in \langle g \rangle \cdot \overline{R}$.
        \item If $l(h,R)=l(g,R) + \mu(R) - r(R)$ then $h \notin \langle g \rangle$. 
    \end{enumerate}
\end{proposition}

\begin{proof}
    For the first item, if $l(h, \frac{R}{\mathfrak{p}_i}) \geq l(g, \frac{R}{\mathfrak{p}_i}) + \mu_i(R)$ for every $i$ then we have that $l(h, \frac{R}{\mathfrak{p}_i}) - l(g, \frac{R}{\mathfrak{p}_i}) \geq \mu_i(R)$. Therefore, we can conclude from Proposition~\ref{prop:mu_vec} that there exists some $a \in R$ such that $l(a, \frac{R}{\mathfrak{p}_i})=l(h, \frac{R}{\mathfrak{p}_i}) - l(g, \frac{R}{\mathfrak{p}_i})$ for every $i$. Thus  $l(ag, \frac{R}{\mathfrak{p}_i})=l(h, \frac{R}{\mathfrak{p}_i})$, and  the result follows from Lemma~\ref{lem:ord_ideal}.  \\

    For the second item, assume towards contradiction that $h \in \langle g \rangle$ satisfies $l(h,R)=l(g,R) + \mu(R) - r(R)$. Therefore there exists some $a \in R$ such that $h=ga$. Thus, from Lemma~\ref{lem:semigroup} we can conclude that $l(g,R)+l(a,R)=l(h,R)=l(g,R) + \mu(R) - r(R)$. This gives us that $l(a,R)=\mu(R) - r(R)$ which is impossible by Remark~\ref{rem:mu-r(R)-gamma}. 
\end{proof}

% We end this section by showing how we can relate $\edim(R)$ with the invariant $\Gamma(R)$. 

% We now turn to defining a notion of equivalence between one dimensional rings based upon their semigroup. 

% \begin{definition}
%     We say that $R_1$ and $R_2$ are \textbf{equisingular}, and we denote $R_1 \equiv R_2$, if $\nu(R_1)=\nu(R_2)$. 
% \end{definition}

% \begin{remark}
%     Inspiration for equisingularity 
% \end{remark}

% \begin{proposition}
% If $R_1$ and $R_2$ are equisingular then
% \begin{enumerate}
%     \item $\mult(R_1)=\mult(R_2)$, 
%     \item $\delta(R_1)=\delta(R_2)$, 
%     \item $r(R_1)=r(R_2)$, 
%     \item $\mu(R_1)=\mu(R_2)$. 
%     \item $\Gamma(R_1)=\Gamma(R_2)$. 
% \end{enumerate}
% \end{proposition}

% \begin{proof}
    
% \end{proof}

\section{One Dimensional Rings of Finite Cohen-Macaulay Type}\label{sec:FCMT1}

In this section we show how we can use the tools and techniques we developed to relate one dimensional rings of finite Cohen-Macaulay type to the classical ADE curve singularities using their semigroups. For more information on finite and countable Cohen-Macaulay type, see~\cite{svoray2025ade, yoshino1990maximal, leuschke2012cohen}. \\

Throughout this section we assume that $(R, \mathfrak{m}, \kappa)$ is a one dimensional Cohen-Macaulay local ring with $\chara(\kappa) \neq 2$ (see Remark~\ref{rem:char_nu} for the reason we add this assumption). By the fifth item of Proposition~\ref{prop:CMtype_onedim} we can also assume that $R$ is complete. 

\begin{definition}\label{def:cm_type}
     Denote by $\mathcal{MCM}(R)$ the set of isomorphism classes $[M]$ of maximal indecomposable Cohen–Macaulay modules over $R$. If $|\mathcal{MCM}(R)|$ is finite we say that $R$ has finite Cohen-Macaulay type.  
\end{definition}

The following proposition summarize the main results of Chapter $4$ of~\cite{leuschke2012cohen} (using our notations). 

\begin{proposition}\label{prop:CMtype_onedim}
Assume that $R$ is Cohen-Macaulay. Then: 
    \begin{enumerate}
        \item If $R$ does not have finite Cohen-Macaulay type then $|\mathcal{MCM}(R)| \geq|\kappa|$.  
        \item $R$ has finite Cohen-Macaulay type if and only if $\mult(R) \leq 3$ and $\frac{\mathfrak{m}\overline{R}+R}{R}$ is a cyclic $R-$module. 
        \item If $R$ has finite Cohen-Macaulay type then $R$ is analytically unramified (and therefore reduced). 
        \item If $\mult(R)=2$ and $R$ is analytically unramified then $R$ has finite Cohen-Macaulay type. 
        \item $R$ has finite Cohen-Macaulay type if and only if the completion of $R$ does. 
    \end{enumerate}
\end{proposition}

\begin{remark}
\textup{From the first item of Proposition~\ref{prop:CMtype_onedim} we have that if $R$ is a one dimensional ring then either $R$ has finite Cohen-Macaulay type or $|\mathcal{MCM}(R)|\geq|\kappa|$. This in fact tells us that in the one dimensional case, having a sparse Cohen-Macaulay type (as in Definition 2.3 in~\cite{svoray2025ade}) is the same as having Finite Cohen-Macaulay type.}
\end{remark}

In order to understand one dimensional rings of finite Cohen-Macaulay type we define a notation that allows use to compare rings based upon their semigroups:

\begin{definition}\label{def:equidominated}
    We say that $R_1$ and $R_2$ are equisingular if $\nu(R_2) = \nu(R_1)$.  
\end{definition}

\begin{remark}
\begin{enumerate}
    \item \textup{The name "equisingular" comes from a similar notion of "equisingularity of algebraic curves", in which we say that two irreducible power series $f,g \in \mathbb{C}[[x,y]]$ over the complex numbers are equisingular if $\Gamma\left(\frac{\mathbb{C}[[x,y]]}{\langle f \rangle}\right) = \Gamma\left(\frac{\mathbb{C}[[x,y]]}{\langle g \rangle}\right)$. This notion over algebraic curves allows us to understand the topological type of the curves, as studied in~\cite{zariski2006moduli, zariski1965studies}) and to study deformations of plane curve singularities, as in Section 2 of Chapter 2 in~\cite{greuel2007introduction}. A notion of equisingularity exists for higher dimension and was developed by Zariski to study resolution of singularities (for more information, see~\cite{zariski1971some, parusinski2021algebro, lipman2000equisingularity}).  }
    \item \textup{The reason we look at $\nu(R)$ and not at $\Gamma(R)$ in the definition of equisingularity is because $\Gamma(R)$ does not distinguish between rings based upon the number of minimal primes (for example, Remark~\ref{rem:A_12groups} tells us that $A_1$ and $A_2$ singularities have the same $\Gamma(R)$ semigroup). In addition,  Proposition~\ref{prop:nu_beasic_prop} tells us that elements in $\nu(R)$ corresponds with elements in $R$ (a property we discuss in Remark~\ref{rem:ADE_semigroup}). }
\end{enumerate}
\end{remark}

 %The second is "birationally dominates", in which we say that a one dimensional ring $R_1$ birationally dominates a one dimenisonal ring $R_2$ if $R_2 \subset R_1 \subset \overline{R_2}$, which is used in~\cite{greuel1984einfache} to classify complete one dimensional local rings of finite Cohen-Macaulay type that contain an algebraically closed field of characteristic zero (in fact, as we see in Remark~\ref{rem:ADE_semigroup}, if $R$ containes a field, then these two notions agree).

The following lemma shows that being equisingular tells us that we have a "similar" singularity, in an algebraic sense. 

\begin{lemma}
    If $R_1$ and $R_2$ are equisingular $R_1$ then 
\begin{enumerate}
    \item $\mult(R_1) =  \mult(R_2)$, 
    \item $\delta(R_1) = \delta(R_2)$, 
    \item $r(R_1)=r(R_2)$, 
    \item $\mu(R_1) = \mu(R_2)$. 
    \item $\Gamma(R_1) = \Gamma(R_2)$. 
\end{enumerate}
\end{lemma}

\begin{proof}
    These follow directly from the definition of $\nu(R)$ together with Lemma~\ref{lem:semigroup}, Proposition~\ref{prop:nu_beasic_prop}, Proposition~\ref{prop:mu_vec}, and Corollary~\ref{cor:Gamma_cond}. 
\end{proof}

We now define a family of special one dimensional rings that we wish to relate to one dimensional rings of finite Cohen-Macaulay type. Note that two of these types of singularities is $A_k$, exactly as we defined in Definition~\ref{def:A_k}. In addition, we added their of $\mu(R)$ and $r(R)$, as these values are easy to compute (in a similar fashion to the computation done in the second item of Proposition~\ref{cor:delta1}). 

\begin{definition}\label{def:ADE_curves}
    We say that $R$ is an ADE type singularity if $\hat{R} \cong \frac{S}{\langle f \rangle}$ where $(S, \mathfrak{n}, \kappa)$ be a $2-$dimensional complete regular local ring and $\mathfrak{n} = \langle x,y \rangle$ such that $f$ equals to one of the following below:
\begin{center}
\hspace*{-0.8cm}
    \begin{tabular}{ | l | l | l | l | l |}
    \hline
    $ADE$  & Formula & $\mu(R)$ & $r(R)$ &  Note  \\ \hline \hline
    $A_{2n}$ & $x^2  + y^{2n+1}$ & $2n$ & $1$ & $1 \leq n$  \\ \hline
    $A_{2n+1}$ & $x^2  + y^{2n+2}$ & $2n+1$ & $2$ & $0 \leq n$  \\ \hline
    $D_{2n}$ & $x\left(y^2  + x^{2n-2}\right)$ & $2n$ & $3$ & $2 \leq n$  \\ \hline
    $D_{2n+1}$ & $x\left(y^2  + x^{2n-1}\right)$ & $2n+1$ & $2$ & $2 \leq n$  \\ \hline
    $E_6$ & $x^3 + y^{4}$ & $6$ & $1$ &    \\ \hline
    $E_7$ & $x\left(x^2 +  y^{3}\right)$ & $7$ & $2$ &  \\ \hline
    $E_8$ & $x^3 + y^{5}$ & $8$ & $1$ &   \\ \hline
    \end{tabular}

    \vspace{5mm}
\end{center}

\end{definition}

\begin{remark}\label{rem:char_nu}
\begin{enumerate}
    \item \textup{In~\cite{svoray2025ade} it was shown that in fact, if $R$ is a one dimensional ring of finite Cohen-Macaulay type with $\edim(R)=2$ (assuming $\chara(\kappa) \neq 2,3,5$) then $R$ is an ADE singularity.  If $\chara(\kappa)$ is either $3$ or $5$, then there are additional ADE singularities that come up. }
    \item \textup{As mentioned in the previous item, when discussing ADE singularities beyond the characteristic zero case, such as in~\cite{svoray2025ade,greuel1988simple}, there are additional forms that are considered ADE singularities (specifically, there additional $E_6, E_7$, and $E_8$ singularities in characteristics $3$ and $5$, in addition to a different classification in characteristic $2$). Yet, as we see below, we need not look at these additional singularities since similar computations to those preformed in Lemma~\ref{lem:ADE_nu_grp} show us that the have the same semigroup $\nu(R)$.   }
    \item \textup{The reason we assume that $\chara(\kappa) \neq 2$ is that in the characteristic $2$ case some of the ADE singularities behave differently. Specifically, $A_{2n+1}$ and $D_{2n}$ singularities are not reduced if $R$ is equicharacteristic and irreducible if $R$ is of mixed characteristic $(0,2)$.}
\end{enumerate}
\end{remark}

Therefore the main result of this section is that one dimensional rings of finite Cohen-Macaulay type is equisingular to an ADE type singularity:

\begin{theorem}\label{thm:ADE_curve}
    If $R$ has finite Cohen-Macaulay type then $R$ is equisingular to a ring of ADE type singularity.
\end{theorem}




% \begin{notation}\label{notation:semgro}
    
% \end{notation}

% \begin{remark}
%     \textup{Notation~\ref{notation:semgro} makes sense since in fact, every sub-semigroup of $\mathbb{N}$ is finitely generated, and a sub-semigroup of $\mathbb{N}$ is numerical if and only if it is generated by a finite sequence $a_1, \dots, a_r$ such that $\gcd(a_1, \dots, a_r)=1$. For proofs and more information, . }
% \end{remark}

The following lemma tells us the semigroups of rings of ADE type. 

\begin{lemma}\label{lem:ADE_nu_grp}
    \begin{enumerate}
        \item If $R$ is of $A_{2n}$ type then $\nu(R)= \lceil 2, 2n+1 \rfloor$.
        \item If $R$ is of $A_{2n+1}$ type then $\nu(R) = \lceil (1,1)\rfloor \cup ((n+1, n+1) + \mathbb{N}^2)$ and $\Gamma(R)=\lceil 2, 2n+3 \rfloor$. 
        \item If $R$ is of $D_{2n}$ type then $\nu(R) = \{(a,b,b) \colon b < n-2, a \leq b(n-1)\} \cup \{(1, n-1, n-1)\} \cup  ((2,n-1,n-1) + \mathbb{N}^3)$ and $\Gamma(R)=\{3\} \cup (5 + \mathbb{N})$.
        \item If $R$ is of $D_{2n+1}$ type then $\nu(R) = \{(r,2s) \colon r \geq 2, s \geq 1 \} \cup \{(1,2s) \colon 2s \leq 2n-1\} \cup \{(1, 2n-1)\} \cup ((2, 2n)+\mathbb{N}^2)$ and $\Gamma(R)=\{3\} \cup (5 + \mathbb{N})$.
        \item If $R$ is of $E_{6}$ type then $\nu(R) = \lceil 3,4 \rfloor$
        \item If $R$ is of $E_{7}$ type then $\nu(R) = \lceil (1,2) \rfloor \cup \{(k,3) \colon k >1\} \cup ((3,5) + \mathbb{N}^2)$ and $\Gamma(R) = \{0,3\} \cup (5+ \mathbb{N})$. 
        \item If $R$ is of $E_{8}$ type then $\nu(R) = \lceil 3,5 \rfloor$
    \end{enumerate}
\end{lemma}

\begin{proof}
    \begin{enumerate}
        \item Since $\mult(R)=2$ then by Lemma~\ref{rem:semigroup} we have that $2 \in \Gamma(R)$ and its the non zero minimum of $\Gamma(R)$. In addition, by Proposition~\ref{cor:delta1} we have that $2\delta(R)=2n$, and so by Corollary~\ref{cor:Gamma_cond} we have that $N \in \Gamma(R)$ for every $N \geq 2n$ but $2n-1 \notin \Gamma(R)$. Therefore we must have that $\Gamma(R) = \{2, 4, \dots, 2n\} \cup ((2n+1) + \mathbb{N})$, as if there exists some odd number $2(n-m)+1 \in \Gamma(R)$ that is smaller than $2n$, we would have that $2(n-m)+1 + 2(m-1) = 2n-1 \in \Gamma(R)$ which is a contradiction. 
        
        \item Let $\mathfrak{p}_1 = \langle x+iy^{n+1}\rangle$ and $\mathfrak{p}_2 = \langle x-iy^{n+1}\rangle$ be the minimal primes of $R$ (where $i \in R$ satisfies $i^2=-1$).  Since in this case we have that $r(R)=2$ and $\mult(R)=2$ then by Lemma~\ref{lem:mult_sum_prim} we must have that $\mult\left(\frac{R}{\mathfrak{p}_1}\right)=\mult\left(\frac{R}{\mathfrak{p}_2}\right)=1$, and so $(1,1) \in \nu(R)$. In addition, by Lemma~\ref{cor:delta1} we have that $\vec{\mu}(R)= (n+1, n+1)$, and so by Proposition~\ref{prop:mu_vec} we can conclude that $(n+1, n+1) + \mathbb{N}^2 \subset \nu(R)$. Yet, from Remark~\ref{rem:mu-r(R)-gamma} we get that $\Delta_1(k,k)$ and $\Delta_2(k,k)$ do not intersect $\nu(R)$ for every $0 \leq k \leq n$, and so the result follows, together with Proposition~\ref{prop:nu_beasic_prop}.
        
        \item Denote $\mathfrak{p}_1 = \langle x \rangle$, $\mathfrak{p}_2 = \langle x+iy^{n-1} \rangle$, and $\mathfrak{p}_2 = \langle x+iy^{n-1} \rangle$ the minimal primes of $R$ (where $i \in R$ satisfies $i^2=-1$, which exists since $R$ is strictly Henselian with $\chara(\kappa) \neq 2$). Then we have that $\mult\left(\frac{R}{\mathfrak{p}_1}\right)=\mult\left(\frac{R}{\mathfrak{p}_2}\right)=\mult\left(\frac{R}{\mathfrak{p}_3}\right)=1$, and so $(1,1,1) \in \nu(R)$. In addition, we have that $\vec{\mu}(R)=(2,n,n)$, and so by Proposition~\ref{prop:mu_vec} we can conclude that $(2,n,n) + \mathbb{N}^3 \subset \nu(R)$.  Now, since $\frac{R}{\mathfrak{p}_2 \cap \mathfrak{p}_3}$ is an $A_{2(n-1)}$ singularity, then by Proposition~\ref{prop:nu_beasic_prop} and by the previous item, we must have that if $(a,b,c) \in \nu(R)$ satisfies $b,c  \leq n-1$ then $b=c$. In addition, for $b < n-2$ and $1 \leq a \leq b(n-1)$ we have that $\nu(y^a +x^b) = (a, b, b) \in \nu(R)$, and we also have $\nu(x+y^{n-1})=(1, n-1, n-1) \in \nu(R)$. Since $(0,0,0), (1,1,1) \in \nu(R)$ then by Remark~\ref{rem:mu-r(R)-gamma} we have that $\Delta_i(1, n-1, n-1)$ and $\Delta_i(0, n-2, n-2)$ do not intersect $\nu(R)$ for $i=1,2,3$. 
        
        \item Denote  $\mathfrak{p}_1 = \langle x \rangle$ and $\mathfrak{p}_2 = \langle y^2 + x^{2n-1} \rangle$, the minimal primes of $R$. Then we have that $\left(\mult\left(\frac{R}{\mathfrak{p}_1}\right), \mult\left(\frac{R}{\mathfrak{p}_2}\right)\right)=(1,2) \in \nu(R)$ is the minimal non-zero element and $\vec{\mu}(R)=(2, 2n)$ is the minimal element for which $\vec{\mu}(R)+ \mathbb{N}^2 \subset \nu(R)$. Since $(0,0) \in \nu(R)$ then by Remark~\ref{rem:mu-r(R)-gamma} we must have that $\Delta_1(1, 2n-1))$ and  $\Delta_2(1, 2n-1)$ do not intersect $\nu(R)$, with $\nu(y)=(1, 2n-1) \in \nu(R)$.  In addition, for every $l$ and $r$ such that $r \geq 2$ and $2l \leq r(2n-1)$ we have that $\nu(y^r+x^l)=(r, 2l) \in \nu(R)$, and for every $l$ such that $2l < 2n-1$ we have that $\nu(y+x^l)=(1,2l) \in \nu(R)$. Finally, since $2l+1 \notin \Gamma\left(\frac{R}{\mathfrak{p}_2}\right)$ for every $2l+1 \leq 2n-2$ (as $\frac{R}{\mathfrak{p}_2}$ is an $A_{2n-2}$ type singularity) then by Proposition~\ref{prop:nu_beasic_prop} we must have that $(r, 2l+1) \notin \nu(R)$ for every $r$.
        
        \item Since $\mult(R)=3$ then by Lemma~\ref{rem:semigroup} we have that $3 \in \Gamma(R)$ is the minimal non-zero element in $\Gamma(R)$ and since $\delta(R)=3$ then by Corollary~\ref{cor:Gamma_cond} we have that $N \in \Gamma(R)$ for every $N \geq 6$ but $5 \notin \Gamma(R)$. By Remark~\ref{rem:mu-r(R)-gamma} we can conclude that $|\mathbb{N} \setminus \Gamma(R)|=3$, and so we must have that $4 \in \Gamma(R)$, and the result follows. 
        
        \item Denote $\mathfrak{p}_1 = \langle x \rangle$ and $\mathfrak{p}_2 = \langle x^2 + y^3 \rangle$, the minimal primes of $R$. Then $\mult\left(\frac{R}{\mathfrak{p}_1}\right)=1$ and $\mult\left(\frac{R}{\mathfrak{p}_2}\right)=2$, and so by Remark~\ref{cor:sum_length} we have that $(1,2) \in \nu(R)$. In addition, since $i(\mathfrak{p}_1, \mathfrak{p}_2) = 3$ and $\mu\left(\frac{R}{\mathfrak{p}_2}\right)=2$, we can conclude that $\vec{\mu}(R)=(3,5)$. Since $(0,0), (1,2), (2,4) \in \nu(R)$, then by Remark~\ref{rem:mu-r(R)-gamma} we must have that the sets $\Delta_i(1,2)$, $\Delta_i(2,4)$, and $\Delta_i(0,0)$ do not intersect $\nu(R)$ for $i=1,2$. Finally, for every $k>1$ we have that $i(\langle x \rangle, \langle y^k +x \rangle) = k$ and $i(\langle x^2+y^3 \rangle, \langle y^k +x \rangle) = 3$, and so $(k,3) \in \nu(R)$. 
        
        \item As before, since $\mult(R)=3$ then by Lemma~\ref{rem:semigroup} we have that $3 \in \Gamma(R)$ is the minimal non-zero element in $\Gamma(R)$ and since $\delta(R)=4$ then by Corollary~\ref{cor:Gamma_cond} we have that $N \in \Gamma(R)$ for every $N \geq 8$ but $7 \notin \Gamma(R)$. Note that $4 \notin \Gamma(R)$ as then we would have that $3+4=7 \in \Gamma(R)$. Since by Remark~\ref{rem:mu-r(R)-gamma} we have that $|\mathbb{N} \setminus \Gamma(R)|=4$, we must have that $5,6 \in \Gamma(R)$, and the result follows.  
    \end{enumerate}
\end{proof}


\begin{figure}[ht]
\centering

\begin{tikzpicture}[scale=0.9]

    % Number line
    \draw (0,0) -- (11,0);

    % Initial numeric ticks
    \foreach \x in {0,1,2,3,4,5} {
        \draw (\x,0.1) -- (\x,-0.1)
              node[below=3pt] {\scriptsize $\x$};
    }

    % Symbolic tick at 2N
    \draw (7,0.1) -- (7,-0.1)
          node[below=3pt] {\scriptsize $2n$};

    \draw (8,0.1) -- (8,-0.1)
          node[below=3pt] {\scriptsize $2n+1$};

    % Uniform ticks after 2N (unlabeled)
    \foreach \x in {8,9,10} {
        \draw (\x,0.1) -- (\x,-0.1);
    }

    % Red dots: even numbers up to 2N
    \fill[red] (2,0) circle (2pt);
    \fill[red] (4,0) circle (2pt);
    \fill[red] (7,0) circle (2pt); % 2N
    \fill[red] (0,0) circle (2pt);


    % Ellipsis between 4 and 2N
    \node at (6,0.35) {\scriptsize $\cdots$};

    % Red dots: integers >= 2N+1 (uniform spacing)
    \fill[red] (8,0) circle (2pt);
    \fill[red] (9,0) circle (2pt);
    \fill[red] (10,0) circle (2pt);

    % Arrow indicating continuation
    \draw[black,->] (10.3,0) -- (11,0);

  

\end{tikzpicture}

\caption{The semigroup $\Gamma(R)$ of an $A_{2n}$ singularity}
    \end{figure}
    
    
    
    \begin{figure}[ht]
    \centering
    
    \begin{tikzpicture}[scale=0.8]
    
        % --- Axes ---
        \draw[->] (0,0) -- (9,0) node[right] {};
        \draw[->] (0,0) -- (0,9) node[above] {};
    
        % --- Symbolic ticks: only relevant for the set ---
        \foreach \i in {1,2} {  % 1..N for illustration; N=4 here symbolically
            \draw (\i,0.1) -- (\i,-0.1) node[below=2pt] {\scriptsize $\i$};
            \draw (0.1,\i) -- (-0.1,\i) node[left=2pt] {\scriptsize $\i$};
        }
    
        \draw (3,0.1)  (3,-0.1) node[below=2pt] {\scriptsize $\cdots$};
        \draw (0,3)  (-0.1,3) node[left=2pt] {\scriptsize $\vdots$};
    
        
        
        % Tick for N+1 (start of quadrant)
        \draw (5,0.05) -- (5,-0.05) node[below=2pt] {\scriptsize $n+1$};
        \draw (0.1,5) -- (-0.1,5) node[left=2pt] {\scriptsize $n+1$};
    
        \draw (4,0.1) -- (4,-0.1) node[below=2pt] {\scriptsize $n$};
        \draw (0.1,4) -- (-0.1,4) node[left=2pt] {\scriptsize $n$};
    
        \draw (6,0.1) (6,-0.1) node[below=2pt] {\scriptsize $\cdots$};
        \draw (0.1,6)(-0.1,6) node[left=2pt] {\scriptsize $\vdots$};
    
        % --- Diagonal points (1,1) to (N,N) ---
        \foreach \i in {0,1,2,4} {  % 1..N
            \fill[red] (\i,\i) circle (2pt);
        }
    
        % Ellipsis along diagonal
        \node at (3,3) {\scriptsize $\iddots$};
    
        \foreach \i in {5,6,7,8} {
            \foreach \j in {5,6,7,8} {
                \fill[red] (\i,\j) circle (2pt);
            }
        }
    
        
    
        % Arrow indicating quadrant continues
    
    
    
    \end{tikzpicture}
    
    \caption{The semigroup $\nu(R)$ of an $A_{2n+1}$ singularity.}
    \end{figure}
    
    
    \begin{figure}[ht]
\centering

\begin{tikzpicture}[scale=0.8]

    % --- Axes ---
    \draw[->] (0,0) -- (11,0) node[right] {};
    \draw[->] (0,0) -- (0,11) node[above] {};

    % --- Symbolic ticks: only relevant for the set ---
    \foreach \i in {1,2,3,4} {  % 1..N for illustration; N=4 here symbolically
        \draw (0.1,\i) -- (-0.1,\i) node[left=2pt] {\scriptsize $\i$};
    }


    \draw (1,0.1) -- (1,-0.1) node[below=2pt] {\scriptsize $1$};
    \draw (2,0.1) -- (2,-0.1) node[below=2pt] {\scriptsize $2$};

    % \foreach \i in {4,5, 6, 7, 8, 9} {  % 1..N for illustration; N=4 here symbolically
    %     \draw (\i, 0.1) -- (\i, -0.1) node[left=2pt] {\scriptsize };
    % }

    % \draw (10, 0.1) (10, -0.1) node[left=2pt] {\scriptsize $\cdots$};
    
    \draw (5,0.1)  (3,-0.1) node[below=2pt] {\scriptsize $\cdots$};
    \draw (0,5)  (-0.1,3) node[left=2pt] {\scriptsize $\vdots$};

    
    
    % Tick for N+1 (start of quadrant)
    \draw (0.1,5.5)  (-0.1,5.5) node[left=2pt] {\scriptsize $\vdots$};
    \draw (0.1,7) -- (-0.1,7) node[left=2pt] {\scriptsize $2n-2$};
    \draw (0.1,8) -- (-0.1,8) node[left=2pt] {\scriptsize $2n-1$};
    \draw (0.1,9) -- (-0.1,9) node[left=2pt] {\scriptsize $2n$};
    \draw (0.1,10)  (-0.1,10) node[left=2pt] {\scriptsize $\vdots$};


    \foreach \i in {7,8}{
        % \draw (\i,0.1) -- (\i,-0.1) node[below=2pt] { };
        \draw (0.1,\i) -- (-0.1,\i) node[left=2pt] { };
        }

    % --- Diagonal points (1,1) to (N,N) ---
    \foreach \i in {2,4, 7} { 
        \foreach \j in {1,2,3,4,5,6,7,8,9,10}{
             \fill[red] (\j,\i) circle (2pt);
            }
    }
    \fill[red] (0,0) circle (2pt);
    \fill[red] (1,8) circle (2pt);

    % Ellipsis along diagonal

    \foreach \i in {2,3,4,5,6,7,8,9,10} {
        \foreach \j in {9,10} {
            \fill[red] (\i,\j) circle (2pt);
        }
    }

    

    % Arrow indicating quadrant continues



\end{tikzpicture}

\caption{The semigroup $\nu(R)$ of an $D_{2n+1}$ singularity.}
\end{figure}



\begin{figure}[ht]
\centering

% -------- First copy --------
\begin{tikzpicture}[scale=0.9]
    \draw[->] (0,0) -- (11.5,0);

    \foreach \x in {0,1,...,11} {
        \draw (\x,0.1) -- (\x,-0.1) node[below=3pt] {\scriptsize $\x$};
    }

    \foreach \x in {0,3,4,6,7,8,9,10,11} {
        \fill[red] (\x,0) circle (2pt);
    }

    \node at (5.75,-1) {\scriptsize    };
\end{tikzpicture}
\hfill
% -------- Third copy --------
\begin{tikzpicture}[scale=0.9]
    \draw[->] (0,0) -- (11.5,0);

    \foreach \x in {0,1,...,11} {
        \draw (\x,0.1) -- (\x,-0.1) node[below=3pt] {\scriptsize $\x$};
    }

    \foreach \x in {0,3,5,6,8,9,10,11} {
        \fill[red] (\x,0) circle (2pt);
    }

    \node at (5.75,-1) {\scriptsize  };
\end{tikzpicture}

\caption{The semigroup $\Gamma(R)$ of an $E_6$ and an $E_8$ singularity.}
\end{figure}


\begin{figure}[ht]
\centering

\begin{tikzpicture}[scale=0.8]

    % --- Axes ---
    \draw[->] (0,0) -- (9,0) node[right] {};
    \draw[->] (0,0) -- (0,10) node[above] {};

    % --- Symbolic ticks: only relevant for the set ---
    \foreach \i in {1,2,3,4,5,6,7,8} {  % 1..N for illustration; N=4 here symbolically
        \draw (0.1,\i) -- (-0.1,\i) node[left=2pt] {\scriptsize $\i$};
    }

    \foreach \i in {1,2,3} {  % 1..N for illustration; N=4 here symbolically
        \draw (\i,0.1) -- (\i,-0.1) node[below=2pt] {\scriptsize $\i$};
    }

    \draw (0.1,) (-0.1,9) node[left=2pt] {\scriptsize $\vdots$};

    \draw (4,0.1)  (4,-0.1) node[below=2pt] {\scriptsize $\cdots$};

    \foreach \i in {2,3,4,5,6,7,8} {
        \fill[red] (\i,3) circle (2pt);
    } 

    % --- Diagonal points (1,1) to (N,N) ---
    \foreach \i in {0,1,2,4} {  % 1..N
        \fill[red] (\i,2*\i) circle (2pt);
    }


    \foreach \i in {3,4,5,6,7,8} {
        \foreach \j in {5,6,7,8,9} {
            \fill[red] (\i,\j) circle (2pt);
        }
    }

    

    % Arrow indicating quadrant continues



\end{tikzpicture}

\caption{The semigroup $\nu(R)$ of an $E_7$ singularity.}
\end{figure}



We prove Theorem~\ref{thm:ADE_curve} in steps based upon the number of minimal primes and upon the multiplicity of $R$, inspired by Chapter 9 of ~\cite{yoshino1990maximal} and by Chapters 8 and 9 of~\cite{lynch2010class}. Note that by the first item of Proposition~\ref{prop:CMtype_onedim} we have that $\mult(R) \leq 3$ and so from Corollary~\ref{cor:mult_r} we have that $r(R) \in \{1,2,3\}$. We start with the case where $\mult(R)=2$, which we can view as a generalization of Corollary~\ref{cor:A_1_Gamma}:

\begin{proposition}\label{prop:equiA}
    Assume that $R$ is not a DVR. Then the following are equivalent:
    \begin{enumerate}
        \item $\mult(R)=2$, 
        \item $R$ is an $A_k-$singularity for some $k \geq 1$, 
        \item $R$ equisingular to an $A_k$ singularity. 
        \item $\Gamma(R) = \lceil 2, 2\delta(R)+1 \rfloor$. 
    \end{enumerate}
\end{proposition}


\begin{proof}
    Note that the equivalence of $1$ and $2$ follows Proposition~\ref{cor:delta1}. In addition, since $R$ is not a DVR then $1 \notin \Gamma(R)$ from Corollary~\ref{Cor:DVR_Gamma}. Therefore, from Lemma~\ref{lem:semigroup} we have that $2 \in \Gamma(R)$ if and only if $\mult(R)=2$. Now, assume that $\mult(R)=2$.  By Corollary~\ref{cor:mult_r} we have that either $r(R)=1$ or $r(R)=2$. We view each case separately:  \\
    
    First, if $r(R)=1$ then $\overline{R}$ is a DVR, and since $\mult(R)=2$ we have that $\pi^2 \in R$ for some uniformizer $\pi$ of $\overline{R}$. In addition, since $\pi^{2\delta(R)} \overline{R} = \mathfrak{c}_R \overline{R}$, with $\mathfrak{c}_R = (R \colon \overline{R})$, we can conclude that $\pi^{2\delta(R)+1} = \pi \cdot \pi^{2\delta(R)}  \in \mathfrak{c}_R \subset R$, and so $\pi^{2\delta(R)+1} \in R$. Therefore from Corollary~\ref{prop:mu_vec} and Corollary~\ref{cor:Gamma_cond} we must have that that $\lceil 2, 2\delta(R)+1 \rfloor = \Gamma(R)$.  \\
    
    Second, if $r(R)=2$ then by Remark~\ref{rem:techincal_ring} we can write $\overline{R} = \overline{\frac{R}{\mathfrak{p}_1}} \times \overline{\frac{R}{\mathfrak{p}_2}}$ where $\mathfrak{p}_1$ and $\mathfrak{p}_2$ are the minimal primes of $R$. Since $\mult(R)=2$ then from Lemma~\ref{lem:mult_sum_prim} we can conclude that $\mult\left(\frac{R}{\mathfrak{p}_1}\right) = \mult\left(\frac{R}{\mathfrak{p}_2}\right)=1$. Therefore, by Lemma~\ref{lem:semigroup} and Remark~\ref{cor:sum_length} we have that  $(1,1)=\left(\mult\left(\frac{R}{\mathfrak{p}_1}\right), \mult\left(\frac{R}{\mathfrak{p}_2}\right)\right) \in \nu(R)$, and so from Proposition~\ref{prop:nu_beasic_prop} we can find some uniformizers $\pi_1$ and $\pi_2$ of $\frac{R}{\mathfrak{p}_1}$ and $\frac{R}{\mathfrak{p}_2}$ (respectively) such that $(\pi_1, \pi_2) \in R \subset \overline{R}$. In addition, since both $\frac{R}{\mathfrak{p}_1}$ and $\frac{R}{\mathfrak{p}_2}$ are DVRs, then by Proposition~\ref{thm:branch} we have that $\delta(R)=i(\mathfrak{p}_1, \mathfrak{p}_2)$. Therefore $\vec{\mu}(R)=(\delta(R), \delta(R))$ and we get that $(\pi_1^{\delta(R)}, \pi_2^{\delta(R)}) \in \mathfrak{c}_R \subset R$. By Proposition~\ref{prop:mu_vec} we have that $(\delta(R),\delta(R)) + \mathbb{N}^2$ is contained in $\nu(R)$ and so we can conclude  that $\lceil (1,1)\rfloor + ((n+1, n+1) + \mathbb{N}^2) \subset \nu(R)$, with equality following from applying Remark~\ref{rem:mu-r(R)-gamma} to every element of the form $(k,k)$ for $k \leq n$.  
\end{proof}

\begin{remark}
    \textup{Using a similar proof to the one in Proposition~\ref{prop:equiA}, we can show that $R$ is an ordinary multiple point if and only if $\nu(R) = \lceil (1, \dots, 1)\rfloor \cup ((r(R)-1, \dots, r(R)-1) + \mathbb{N}^{r(R)})$. This is true since by Remark~\ref{rem:omp} we have that $R$ is an ordinary multiple point if and only if $\frac{R}{\mathfrak{p}_i}$ is a DVR for every $i$ and that $i(\mathfrak{p}_i, \mathfrak{p}_j)=1$ for every $i \neq j$, which by Proposition~\ref{prop:nu_beasic_prop} and Proposition~\ref{prop:mu_vec} is equivalent to having $(1,\dots,1)$ be the minimal element in $\nu(R)$ with $\vec{\mu}(R)=(r(R)-1, \dots, r(R)-1)$, and the equivalence follows from applying Remark~\ref{rem:mu-r(R)-gamma} to each $(k, \dots,k) \in \nu(R)$ for $k \leq r(R)-1$. }
\end{remark}

We now turn to the case $\mult(R)=3$, which we split into three lemmata based upon the three cases of $r(R) \in \{1,2,3\}$: 

\begin{lemma}[$E_6$ and $E_8$ cases]\label{lem:equi31}
    If $R$ has finite Cohen-Macaulay type such that $r(R)=1$ and $\mult(R)=3$ then $R$ is equisingular to either an $E_6$ or to an $E_8$ singularity. 
\end{lemma}

\begin{proof}
    Since $\mult(R)=3$ then, as before, we can find a uniformizer $\pi$ of $\overline{R}$ such that $\pi^3 \in R$. In addition, since $\mult(R)=3$ then (as in Lemma~\ref{cor:edim_bound}) we have that $\frac{\mathfrak{m}\overline{R}}{\mathfrak{m}^2 \overline{R}}$ is a three dimensional $\kappa-$vector space, and so we can find some $\alpha$ and $\beta$ in $\mathfrak{m}\overline{R}$ such that $\pi^3, \alpha, \beta$ form a linear basis for $\frac{\mathfrak{m}\overline{R}}{\mathfrak{m}^2 \overline{R}}$. In addition, since $\mult(R)=3$ then we have that $\pi^6 \in \mathfrak{m}^2 \overline{R}$, and since $R$ is of finite Cohen-Macaulay type, then by the first item of Proposition~\ref{prop:CMtype_onedim} we have that $\frac{\mathfrak{m}\overline{R}+R}{R}$ is a cyclic $R-$module, which is equivalent to $\frac{\mathfrak{m}\overline{R}+R}{\mathfrak{m}^2 \overline{R}+R}$ being a one dimensional $\kappa-$vector space. Therefore, we must have that either $\alpha \in R$ or $\beta \in R$. In addition, as $\mult(R)=3$, we can conclude that $\pi^3 \in \mathfrak{m} \overline{R}$ and so $\langle \pi^6 \rangle \subset \mathfrak{m}^2 \overline{R}$. Yet, since $\pi^3, \alpha, \beta$ form a linear basis for $\frac{\mathfrak{m}\overline{R}}{\mathfrak{m}^2 \overline{R}}$ then we must have that the order of $\alpha$ and $\beta$ (as elements over $\overline{R}$) must be different, bigger than $3$, and smaller than $6$. Thus we can conclude (up to renaming them) that $\nu(\alpha) =4$ and $\nu(\beta)=5$. So in total, there exists some $\gamma \in R$ such that either $\nu(\gamma)=4$ or $\nu(\gamma)=5$.\\
    
    If $\nu(\gamma)=4$ then we can conclude that $\lceil 3, 4 \rfloor \subseteq \Gamma(R)$. In this case we must have equality (recalling that $\mult(R)=3$) as otherwise we have that $5 \in \Gamma(R)$, and so $3+\mathbb{N} = \nu(R)$ which is impossible since by Corollary~\ref{cor:Gamma_cond} we would have that $3=2\delta(R)$. Therefore, we can conclude that $\lceil 3, 4 \rfloor = \Gamma(R)$  and thus $R$ equisingular to an $E_6$ singularity. \\
    
    If $\nu(\gamma)=5$ then $\lceil 3, 5 \rfloor \subseteq \Gamma(R)$. As in the previous case, we can not have that $7 \in \nu(R)$ as then either $\Gamma(R)=3+\mathbb{N}$ (which is impossible as then $3=2\delta(R)$) or $\Gamma(R) = \{3\} \cup (5+\mathbb{N})$ (which is impossible as then $5=2\delta(R)$). Thus, we must have that $\lceil 3, 5 \rfloor = \Gamma(R)$  and thus $R$ equisingular to an $E_8$ singularity.
\end{proof}

Now, when moving to the cases where $r(R) \in \{2,3\}$, since we deal with zero-divisors and with elements in $\overline{R}$, we extend the semigroup $\nu(R)$ to include them in the following way:

\begin{notation}
    Let $\mathfrak{p}_1, \dots, \mathfrak{p}_{r(r)}$ be the minimal primes of $R$. Since $\overline{R} = \prod_{i=1}^{r(R)} \overline{\frac{R}{\mathfrak{p}_i}}$, we define $\nu \colon \overline{R} \to \mathbb{N}_\infty^{r(R)}$ by $\nu(g)=\left(\ord(g_i, \overline{\frac{R}{\mathfrak{p}_i}})\right)$ (using the notations of Lemma~\ref{lem:l_irred}), where the order of zero is $\infty$.
\end{notation}

\begin{lemma}[$D_{2n+1}$ and $E_7$ cases]\label{lem:equi3,2}
    If $R$ has finite Cohen-Macaulay type such that $r(R)=2$ and $\mult(R)=3$ then $R$ is equisingular to either an $E_7$ singularity or to an $D_{2n+1}$ singularity (for some $k$).
\end{lemma}

\begin{proof}
    Denote the minimal primes of $R$ by $\mathfrak{p}_1$ and $\mathfrak{p}_2$. Since $\mult(R)=3$ and $r(R)=2$ then from Lemma~\ref{lem:mult_sum_prim} can conclude that $\mult\left(\frac{R}{\mathfrak{p}_1}\right)=1$ and $\mult\left(\frac{R}{\mathfrak{p}_2}\right)=2$. Therefore, from Proposition~\ref{prop:nu_beasic_prop} we have that $(1,2) \in \nu(R)$.  We claim that we can find some uniformizer $\pi_2$ of $\overline{\frac{R}{\mathfrak{p}_2}}$ such that either $(0,\pi_2^2) \in R$ or $(0,\pi_2^3) \in R$. Define $H = \{a \in \mathbb{N}_\infty \colon (\infty, a) \in \nu(R)\}$. Note that $H$ is a numerical semigroup since by Proposition~\ref{prop:mu_vec} we have that $(\pi_1^{\mu_1(R)}, \pi_2^{\mu_2(R)})$ generate $\mathfrak{c}_R$ as an ideal of $R$, and so for every $r\geq 1$ we can conclude that $(0,\pi_2^r)(\pi_1^{\mu_1(R)}, \pi_2^{\mu_2(R)}) \in R$, and therefore $(\infty, r+\mu_2(R)) \in \nu(R)$. Assuming otherwise, since  $\nu(\mathfrak{m}\overline{R}) = (1 + \mathbb{N}_\infty) \times (2 + \mathbb{N}_\infty)$, then we can find two linearly independent elements $\overline{x}$ and $\overline{y}$ in the module $\frac{\mathfrak{m}\overline{R}+R}{\mathfrak{m}^2 \overline{R}+R}$ with inverse images $x,y \in R$ for which we have that $\nu(x) = (\infty, 2)$ and  $\nu(y) = (\infty, 3)$, which contradicts the first item of Proposition~\ref{prop:CMtype_onedim}, as $\frac{\mathfrak{m}\overline{R}+R}{\mathfrak{m}^2 \overline{R}+R}$ has $R-$length at most $1$. \\

    Assume that $(0,\pi_2^2) \in R$. Then since $(1,2) \in \nu(R)$, by Proposition~\ref{prop:nu_beasic_prop} we can find a uniformizer $\pi_1$ of $\overline{\frac{R}{\mathfrak{p}_1}}$ and a unit $u_2 \in \overline{\frac{R}{\mathfrak{p}_2}}$ such that $(\pi_1, u_1\pi_2^{2}) \in R$. Now, since $H$ is a numerical semigroup, then by Remark~\ref{rem:semigroup} it has a conductor $c$ that is an even integer. Therefore we can assume that $u_1$ is of the form $u_1 = v_1 \pi_2^2 + v_m \pi_2^m + \cdots + v_c\pi_2^c$ for odd number $m$, some unit $v_1, v_m \in \overline{R}$, and some $v_{m+1}, \dots,v_c \in \overline{R}$ that are either units or zero. Therefore, by subtracting $(0,\pi_2^2)$ from $(\pi_1, u_1\pi_2^{2})$, and noting that $(0, \pi_2^l) \in R$ for every $l \geq c$, we can conclude that $(\pi_1, \pi_2^{m}) \in R$. Since $m$ is odd, write $m=2n+1$. So, we can conclude that for every $r$ and $l$ for which we have that $r \geq 2$ and $2l \leq r(2n-1)$ we have that $(0, \pi_2^2)^r + (\pi_1, \pi_2^{2n+1})^l = (\pi_1^l, \pi_2^{2r} ( 1+\pi_2^{r(2n-1)-2l}) \in R$, and so $(r, 2l) \in \nu(R)$. Since $\nu(R)$ is a semigroup, we have that for every $r \geq 2$ and for every $l \geq 1$ we have that $(r, 2l) \in \nu(R)$. In addition, for every $l$ such that $2l < 2n-1$ we have that $(0, \pi_2^2)^{l} + (\pi_1, \pi_2^{2n+1}) = (\pi_1, \pi_2^{2l}(1+\pi_2^{2n+1-2l}) \in R$ and so $(1,2l) \in \nu(R)$. Since $(0, \pi_2^2), (\pi_1, \pi_2^{2n+1}) \in R$ then for every $r$ we have that $(0, \pi_2^{2(n+r+1)}) \in R$ and so we have that $\lceil 2, 2n+1 \rfloor \subset H$. Therefore, for every  $s \geq 2n+1$ odd and for every $r \geq 1$ and we have that $(\pi_1^r, \pi_2^{(2n+1)r})(0, \pi_2^s)\in R$ and so $(r,s) \in \nu(R)$. Thus we can conclude that $(2, 2n) + \mathbb{N}^2 \subset \nu(R)$, an so by applying Remark~\ref{rem:mu-r(R)-gamma} we can conclude that $R$ is equisingular to an $D_{2n+1}$ singularity. \\

    If $(0,\pi_2^3) \in R$, since $(1,2) \in \nu(R)$, then as in the previous case, we can find a uniformizer $\pi_1$ of $\overline{\frac{R}{\mathfrak{p}_1}}$ and a unit $u_2 \in \overline{\frac{R}{\mathfrak{p}_2}}$ such that $(0,\pi_2^3)+ (\pi_1^k, u_1\pi_2^{2k}) = (\pi_1^k, (1+u_1\pi_2^{2k-3})\pi_2^3) \in R$. Since $k \geq 2$ then $(1+u_1\pi_2^{2k-3})$ is a unit, and so again by Proposition~\ref{prop:nu_beasic_prop} we have that $(k,3) \in \nu(R)$. In addition, since $(0,\pi_2^3), (\pi_1, u_1\pi_2^{2})  \in R$ then we have that $(0, u_1^i\pi_2^{2i+3j}) \in R$ for every $i,j$, and so we can conclude that $5,6,7 \in H$. Therefore, from the fourth item of Remark~\ref{rem:semigroup}  we can conclude that $H$ is a numerical semigroup with $\lceil 3, 5 \rfloor \subset H$. Therefore, for every $r \geq 5$ there exists a unit $u_r$ such that $(0,u_r \pi_2^r) \in R$, and together with the fact that $(\pi_1^k, u_1\pi_2^{2k}) \in R$ for every $k \geq 1$, we can conclude that $(3,5) + \mathbb{N}^2 \subset \nu(R)$, and therefore by Lemma~\ref{lem:ADE_nu_grp} and Remark~\ref{rem:mu-r(R)-gamma} we can conclude that $R$ is equisingular to an $E_7$ singularity. 
\end{proof} 

\begin{lemma}[$D_{2n}$ case]\label{lem:equi33}
    If $R$ has finite Cohen-Macaulay type such that $r(R)=3$ and $\mult(R)=3$ then $R$ is equisingular to a $D_{2n}$ singularity for some $n$. 
\end{lemma}

\begin{proof}
    Denote the minimal primes of $R$ by $\mathfrak{p}_1, \mathfrak{p}_2, \mathfrak{p}_3$. Since $r(R)=\mult(R)=3$, then by Lemma~\ref{lem:mult_sum_prim} can conclude that $\mult\left(\frac{R}{\mathfrak{p}_1}\right)=\mult\left(\frac{R}{\mathfrak{p}_2}\right)=\mult\left(\frac{R}{\mathfrak{p}_3}\right)=1$, and so from Proposition~\ref{prop:nu_beasic_prop} we have that $(1,1,1) \in \nu(R)$ is the minimal element of $\nu(R)$. We claim that one of $(1,1,\infty), (1,\infty,1),$ or $(\infty,1,1)$ belong to $\nu(R)$. Assuming otherwise, we can find some $\alpha$ and $\beta$ in $\overline{R}$ such that $\nu(\alpha) = (1, 1, \infty)$ and $\nu(\beta) = (1, \infty,1)$ such that their projection onto $\frac{\mathfrak{m}\overline{R}+R}{\mathfrak{m}^2 \overline{R}+R}$ are linearly independent, which would contradict the first item of Proposition~\ref{prop:CMtype_onedim}. Therefore, up to reordering the minimal primes of $R$, we can assume that $(\infty,1,1) \in \nu(R)$, and so we can find uniformizers $\pi_1$ and $\pi_2$ of  $\overline{\frac{R}{\mathfrak{p}_1}}$ and $\overline{\frac{R}{\mathfrak{p}_2}}$, respectively, such that $(0,\pi_2, \pi_3) \in R$. As in the proof of the $D_{2n+1}$ case in Lemma~\ref{lem:equi3,2}, we can conclude that $(r,s,s) \in \nu(R)$ for every $r \geq 2$ and $s \geq 1$ and that $(1,l,l) \in \nu(R)$ for every $l<n$. Finally, if we look at $\mathfrak{q}= \mathfrak{p}_1 \cap \mathfrak{q}$ then by Proposition~\ref{prop:nu_beasic_prop} we have that the projection $R \to \frac{R}{\mathfrak{q}}$ correspond with the projection of $\nu(R)$ onto the $(y,z)-$plane. Therefore, since we have that $\vec{\mu}\left(\frac{R}{\mathfrak{q}}\right) + \mathbb{N}^2 \subset \nu\left(\frac{R}{\mathfrak{q}}\right)$, and by Proposition~\ref{prop:equiA} we have that $\vec{\mu}\left(\frac{R}{\mathfrak{q}}\right)=(n,n)$ for some $n\in \mathbb{N}$, we can conclude that for every $(a,b) \geq (n,n)$ there exists some $c \in \mathbb{N}$ such that $(c,a,b) \in \nu(R)$. Yet, as $(1,1,1) , (\infty,1,1) \in \nu(R)$ we can conclude that $(2,n,n) +\mathbb{N}^3 \subset \nu(R)$, and so by Remark~\ref{rem:mu-r(R)-gamma} we have that $R$ is equisingular to an $D_{2n}$ singularity. 
\end{proof}

\begin{remark}\label{rem:ADE_semigroup}
\textup{If we look carefully at the computations of elements in $R$ preformed in Lemma~\ref{lem:ADE_nu_grp}, in Proposition~\ref{prop:equiA}, Lemma~\ref{lem:equi31}, Lemma~\ref{lem:equi3,2}, and Lemma~\ref{lem:equi33}, we can conclude that if $R$ is a one dimensional local ring with finite Cohen-Macaulay type with minimal primes $\mathfrak{p}_1, \dots, \mathfrak{p}_{r(R)}$, then for every $i=1, \dots, r(R)$, we can find a uniformizer $\pi_i$ of $\overline{\frac{R}{\mathfrak{p}_i}}$ such that $R$ contains the pair of elements $x$ and $y$, which are one of the following:
\begin{enumerate}
    \item ($A_{2n}$ case) $x=\pi_1^{2n+1}$ and $y=i\pi_1^{2}$,  
    \item ($A_{2n+1}$ case) $x=(i\pi_1^n,i\pi_2^n)$ and $y=(\pi_1, \pi_2)$, 
    \item ($D_{2n}$ case) $x=(0, \pi_2, \pi_3)$ and $y=(\pi_1,i\pi_2^{n-1},-i\pi_3^{n-1})$, 
    \item ($D_{2n+1}$ case) $x=(0, \pi_2^2)$ and $y=i(\pi_1, \pi_2^{2n-1})$, 
    \item ($E_6$ case) $x=-\pi_1^4$ and $y=\pi_1^3$, 
    \item ($E_7$ case) $x=(0, \pi_2^3)$ and $y=-(\pi_1, \pi_2^2)$, 
    \item ($E_8$ case) $x=-\pi_1^5$ and $y=\pi_1^3$, 
\end{enumerate}
where $i \in R$ satisfies $i^2=-1$. Note that $x$ and $y$ satisfy the relation $f(x,y)=0$, where $f$ is  one of the polynomials defined in Definition~\ref{def:ADE_curves}. Note that these relations are exactly the generators of the ADE in the case where $S=k[[x,y]]$ for an algebraically closed field of characteristic zero. This, in particular gives us that if $R$ is a one dimensional ring of finite Cohen-Macaulay type that contains an algebraically closed field $k$ of characteristic zero, then $R$ in fact birationally dominates an ADE singularity (where $S=k[[x,y]]$), which is the classical classification of Greuel and Kn\"{o}rrer, first proven in~\cite{greuel1984einfache}.}
\end{remark}

% \section{One Dimensional rings of Finite Cohen-Macaulay type}

% In this section we show how we can use the tools and techniques we developed to classify the Cohen-Macaulay type of $R$ as a one dimensional ring, and to see how we can classify one dimensional rings of finite Cohen-Macaulay type (in the case where $\chara(\kappa) \neq 2$). For more information on finite and countable Cohen-Macaulay type, see~\cite{svoray2025ade, yoshino1990maximal, leuschke2012cohen}. 

% \begin{definition}\label{def:cm_type}
%      Denote by $\mathcal{MCM}(R)$ the set of isomorphism classes $[M]$ of maximal indecomposable Cohen–Macaulay modules $M$ over $R$. If $|\mathcal{MCM}(R)|$ is finite we say that $R$ has \textbf{finite Cohen-Macaulay type}.  
% \end{definition}

% The following proposition summaries the main results of Chapter $4$ of~\cite{leuschke2012cohen} (using our notations). 

% \begin{proposition}\label{prop:CMtype_onedim}
% Assume that $(R, \mathfrak{m})$ is a one dimenisonal local ring. Then: 
%     \begin{enumerate}
%         \item $R$ has finite Cohen-Macaulay type if and only if $\mult(R) \leq 3$ and $\frac{\mathfrak{m}\widetilde{R}+R}{R}$ is a cyclic $R-$module. 
%         \item If $R$ has finite Cohen-Macaulay type then $R$ is analytically unramified (and therefore reduced). 
%         \item If $\mult(R)=2$ and $R$ is analytically unramified then $R$ has finite Cohen-Macaulay type. 
%         \item $R$ has finite Cohen-Macaulay type if and only if the completion of $R$ does. 
%         \item If $R$ does not have finite Cohen-Macaulay type then $|\mathcal{MCM}(R)|=|\kappa|$. 
%     \end{enumerate}
% \end{proposition}

% \begin{remark}
%     \textup{Sparse CM type and one dimensional rings. in~\cite{svoray2025ade}}
% \end{remark}

% By the first item of Proposition ~\ref{prop:CMtype_onedim}, in order to classify one dimensional rings $R$ of finite Cohen-Macaulay type, we need to focus rings of multiplicity $3$ (as rings of multiplicity $2$ have been classified in Proposition~\ref{cor:delta1}), which by Remark~\ref{rem:techincal_ring} tells us that $\edim(R) \in \{1,2,3\}$. Recalling that $\edim(R)=1$ if and only if $R$ is a DVR, we can in fact summaries the case where $\edim(R)=2$, based upon the results from~\cite{svoray2025ade}. Note that in the table below we add the corresponding values to numerical invariants we discussed, which are easy to compute directly.  


% \begin{proposition}\label{prop:E6Dk}
%     Let $(S, \mathfrak{n}, \kappa)$ be a $2-$dimensional complete regular local ring such that $\kappa$ is algebraically closed of characteristic $\neq 2$ and let $f \in \mathfrak{m}^2$. If $R=\frac{S}{\langle f \rangle}$ has finite Cohen-Macaulay type then we can find some $x$ and $y$ that generate $\mathfrak{n}$ such that $f$ equals to one of the following below:

% \begin{center}
% \hspace*{-0.8cm}
%     \begin{tabular}{ | l | l | l | l | l |}
%     \hline
%     $ADE$  & Formula & $\mu(R)$ & $r_R$ &  Note  \\ \hline \hline
%     $A_{2k}$ & $x^2  + y^{2k}$ & $2k$ & $2$ & $1 \leq k$  \\ \hline
%     $A_{2k+1}$ & $x^2  + y^{2k+1}$ & $2k+1$ & $1$ & $0 \leq k$  \\ \hline
%     $D_{2k}$ & $x\left(y^2  + x^{2k-2}\right)$ & $2k$ & $3$ & $2 \leq k$  \\ \hline
%     $D_{2k+1}$ & $x\left(y^2  + x^{2k-1}\right)$ & $2k+1$ & $2$ & $2 \leq k$  \\ \hline
%     $E_6$ & $x^3 + y^{4}$ & $6$ & $1$ &    \\ \hline
%     $E_6^1$ & $x^3 + y^{4} + x^2 y^2$ & $6$ & $1$ & $\chara(\kappa) = 3$  \\ \hline
%     $E_7$ & $x\left(x^2 +  y^{3}\right)$ & $7$ & $2$ &  \\ \hline
%     $E_7^1$ & $x\left(x^2 +  y^{3} + x y^2\right)$ & $7$ & $2$  &$\chara(\kappa) =3$  \\ \hline
%     $E_8$ & $x^3 + y^{5}$ & $8$ & $1$ &   \\ \hline
%     $E_8^1$ & $x^3 + y^{5} + x^2 y^3$  & $8$ & $1$ & $\chara(\kappa)=3$ \\ \hline
%     $E_8^2$ & $x^3 + y^{5}  + x^2 y^2$ &  $8$ & $1$ & $\chara(\kappa)=3$  \\ \hline
%     $E_8^1$ & $x^3 + y^{5}+ x y^4$ & $8$ & $1$ & $\chara(\kappa)=5$ \\ \hline 
%     \end{tabular}
%     \vspace{5mm}
% \end{center}

% Otherwise, $f \in \langle x,y^2\rangle^3$ for some $x,y \in \mathfrak{n}$, in which case $\mu(R)>8$.
% \end{proposition}





% \begin{definition}
%     Let $R$ be a complete local ring such that $\chara(\kappa) \neq 2$. We say that $R$ birationally dominates an ADE singularity if there exists some complete regular local ring $(S, \mathfrak{n})$ and some $f$ as in the table of Proposition~\ref{prop:E6Dk} (with respect to some $x,y$ that generate $\mathfrak{n}$) such that the ring $R_1=\frac{S}{\langle f \rangle}$ satisfies $R_1 \subset R \subset \widetilde{R_1}$. 
% \end{definition}

% Our goal for this section is to prove the following theorem, which classifies all one dimensional Cohen-Macaulay rings of finite Cohen-Macaulay type (recalling that we can assume that the rings are complete by Proposition~\ref{prop:E6Dk}): 

% \begin{theorem}\label{thm:ADE_biration}
%     Let $(R, \mathfrak{m}, \kappa)$ be a one dimensional complete reduced local ring with $\chara(\kappa) \neq 2$. If $R$ has finite Cohen-Macaulay type then $R$ birationaly dominates an ADE singularity. 
% \end{theorem}

% We divide the proof of Theorem~\ref{thm:ADE_biration} into lemmata based upon the value of $r_R$ (recalling that we only need to focus on the case where $\mult(R)=\edim(R)=3$). We set up the following notation that we use throughout this section:

% \begin{notation}
%     \textup{Given some $a_1, \dots, a_n \in \overline{R}$, we denote by $\lceil a_1, \dots, a_n \rfloor$ the subring generated by $a_1, \dots, a_n$, i.e. the smallest subring over $\overline{R}$ containing $a_1, \dots, a_n$.}
% \end{notation}

% \begin{lemma}\label{lem:bir_E76}
%     If $r_R=1$ then $R$ birationally dominates either an $E_6$ or an $E_7$ singularity.
% \end{lemma}

% \begin{proof}
    
% \end{proof}

% \begin{remark}
%     \textup{Similar to Lemma~\ref{lem:bir_E76}, we can show that if $r(R)=1$ and $\mult(R)=2$ then $R$ birationally an $A_n$. Since $r(R)=1$ then $\overline{R}$ is a DVR, and since $\mult(R)=2$ we have that $\pi^2 \in R$. In addition, we have that $\pi^{2\delta(R)} \overline{R} = \mathfrak{c}_R \overline{R}$, with $\mathfrak{c}_R = (R \colon \overline{R})$, we have that $\pi^{2\delta(R)+1} = \pi \cdot \pi^{2\delta(R)}  \in \mathfrak{c}_R \subset R$, and so $\pi^{2\delta(R)+1} \in R$. This gives us that $\lceil \pi^2, \pi^{2\delta(R)+1} \rfloor \subset R$ }
% \end{remark}

% \begin{lemma}
%     If $r_R=2$ then $R$ birationally dominates 
% \end{lemma}




% Since by Lemma~\ref{lem:semigroup} we have that $\Gamma(R)$ is a semigroup, we can look at its generating set and see what information it gives us on $R$. We follow the notations in

% \begin{definition}
%     The minimal set of generators of $\Gamma(R)$ is the finite sequence $(\beta_0, \dots, \beta_h)$ such that:
%     \begin{enumerate}
%         \item $\beta_0=\mult(R)$, 
%         \item For every $k>1$ we have $\beta_k = \min(\Gamma(R) \setminus (\beta_1\mathbb{N} + \cdots + \beta_{k-1}\mathbb{N} ))$, 
%         \item $\Gamma(R) = \beta_1\mathbb{N} + \cdots + \beta_{g}\mathbb{N}$.
%     \end{enumerate}
%      For every $k \leq h$, let $e_k = \gcd(\beta_1, \dots, \beta_k)$ and let $n_k=\frac{e_k}{e_{k-1}}$. 
% \end{definition}

% \begin{remark}\label{rem:beta_ineq}
%     \textup{The fact that $\beta_i$ is a finite sequence follows from.... In addition, note that $e_1 < e_2 < \cdots < e_h$ with $e_{k-1} \beta_k < e_{k} \beta_{k+1}$ for every $k < h$, and so $n_h>1$ is an integer with $n_h \beta_k < \beta_{k+1}$ for every $k<h$. }
% \end{remark}


% \begin{theorem}
%     $\mu(R)=\sum_{k=2}^g (n_k-1)\beta_k -\beta_1 + 1$. 
% \end{theorem}

% \begin{proof}
    
% \end{proof}

% \begin{remark}
%     \textup{Conductor of a semigroup}
% \end{remark}

% \begin{corollary}
%     \begin{enumerate}
%         \item $\mu(R) < n_g\beta_g - \beta_2 - \beta_1+1$.
%         \item $\mu(R)$ is the smallest integer $c \in \mathbb{N}$ such that $c+\gcd(\Gamma(R))\cdot N \in \Gamma(R)$. 
%         \item If $r(R)=1$ then the cardinality of $\mathbb{N} \setminus \Gamma(R)$ is exactly $\mu(R)$. 
%     \end{enumerate}
% \end{corollary}

% \begin{proof}
%     For the first item, by Remark~\ref{rem:beta_ineq} we have that $n_k \beta_k < \beta_{k+1}$ for every $k$, and so we can conclude that 
%     \begin{equation*}
%         \mu(R)=\sum_{k=2}^g (n_k-1)\beta_k -\beta_1 + 1 < \left(\sum_{k=2}^{g-1} \beta_{k+1}-\beta_k\right) -\beta_1 + 1 +(n_g-1)\beta_g,
%     \end{equation*}
%     \noindent and since we have a telescopic series, we can conclude that $\mu(R)< n_g\beta_g-\beta_2-\beta_1+1$.  For the second item,................
% \end{proof}



% We start with $\nu(R)$. Recall that if $\mathfrak{p}_1, \dots, \mathfrak{p}_r$ are the minimal primes of $R$, then the normalization of $R$ is isomorphic to $\prod_{i=1}^r V_i$, where $V_i$ is the integral closure of $\frac{R}{\mathfrak{p}_i}$. Therefore we have an inclusion $R \hookrightarrow \prod_i V_i$, which we denote $g \mapsto (g_1, \dots, g_r)$.  Note that $V_i$ is a DVR with maximal ideal generated by $\pi_i$. 

% \begin{definition}
% Given $g \in R$, we define the order of $g$ to be the tuple $\ord(g)=(\ord_{1}(g_1), \dots, \ord_r(g))$ where $g_i = u_i \pi_i^{\ord_i(g)}$ for some unit $u_i \in V_i$ for every $i$. The order defines a map $\ord \colon R \to \mathbb{N}^r$\footnote{We here denote $\mathbb{N}=\{0, 1, 2, \dots\}$.} and we denote its image by $\nu(R)$. 
% \end{definition}

% \begin{remark}
%     \textup{Note that the uniformizer of a DVR is determined up to a unit, therefore the order is well defined (i.e. does not depend on the choice of uniformizers $\pi_i$). In addition, note that the order satisfies $\ord(gh)=\ord(g)+\ord(h)$ (where addition of tuples is done component-wise) and $\ord(u)=(0, \dots, 0)$ for every unit $u \in R$. }
% \end{remark}

% Recall from Remark~\ref{rem:conductor} that we denote $\mathfrak{c}_R=\Ann_R\left(\frac{\overline{R}}{R}\right)$. Since $\overline{R}$ is normal and one dimensional, it must be a PID, and $\mathfrak{c}_R$, viewed as an ideal in $\overline{R}$, must be generated by an element of the form $(\pi_1^{c_1}, \dots, \pi_r^{c_r})$. We call $\vec{c}=(c_1, \dots, c_r) \in \mathbb{N}^r$ the \textbf{conductor of $R$}. 

% \begin{remark}
%     \textup{If $R$ has a unique minimal prime then by Remark~\ref{rem:conductor} we can conclude that $\vec{c}$, as an integer, equals to $2\delta(R) = \length_R\left(\frac{\overline{R}}{\mathfrak{c}_R} \right)$.}
% \end{remark}

% \begin{proposition}
%     The conductor of $R$ is the minimal element\footnote{By minimal we mean with respect to the induced norm $\mathbb{N}^r  \subset \mathbb{R}^r$.} $\vec{c} \in \nu(R)$ such that $\vec{c} + \vec{v} \in \nu(R)$ for every $\vec{v} \in \mathbb{N}^r$. 
% \end{proposition}

% \begin{proof}
    
% \end{proof}

% \begin{proposition}
    
% \end{proposition}


% \section{Milnor Number of a Curve Hypersurface}

% In this section we study the Milnor number of hypersurface rings, i.e. rings of the form $R=\frac{S}{\langle f \rangle}$ where $(S, \mathfrak{n})$ is a $2-$dimensional complete regular local ring and $f \in \mathfrak{n}$ is reduced.

% \begin{remark}
%     \textup{Observe that in the curve hypersurface case we have that $\mult\left(\frac{S}{\langle f \rangle}\right)$ is in fact the unique constant $N$ such that $f \in \mathfrak{n}^N \setminus \mathfrak{n}^N$. For a proof of this fact, see XXX. For this reason, we simply denote the multiplicity of $\frac{S}{\langle f \rangle}$ by $\mult(f)$.}
% \end{remark}

% \begin{definition}
% Let $(S, \mathfrak{n})$ be a $2-$dimensional complete regular local ring. 
%     \begin{enumerate}
%         \item Given some $f \in \mathfrak{n}$, we denote $\mu(f)=\mu\left(\frac{S}{f}\right)$. 
%         \item Given some $f,g \in \mathfrak{n}$, the intersection multiplicity of $f$ and $g$ is defined to be $i(f,g)=\length_S \left(\frac{S}{\langle f,g \rangle}\right)$.
%     \end{enumerate}
% \end{definition}


% \begin{lemma}
% Let $(S, \mathfrak{n})$ be a regular $2-$dimensional complete local ring and let $f \in \mathfrak{n}$. Denote $R=\frac{S}{\langle f \rangle}$. Then
%     \begin{enumerate}
%         \item $\mu(f)$ is well defined. That is, if $\frac{S}{\langle f \rangle} \cong \frac{S}{\langle g \rangle}$ for some $g \in S$ then $\mu(f)=\mu(g)$.
%         \item $\mu(f)$ is finite and non-negative.
%         \item $\mu(f)=0$ if and only if $\mult(f)=1$. 
%         \item $f=f_1 \cdots f_r$ where $f_i$ is irreducible for every $i$ then 
%         \begin{equation*}
%             \mu(f)=(1-r)+\sum_{i=1}^r \mu(f_i)+2\sum_{ i < j} i(f_i, f_j).
%         \end{equation*}
%         \item $\mu(f) \geq \binom{\mult(f)}{2}$.
%     \end{enumerate}
% \end{lemma}

% \begin{proof}
    
% \end{proof}


% \begin{lemma}
%     Let $(S, \mathfrak{n})$ be a regular $2-$dimensional complete local ring and let $f,g \in \mathfrak{n}$. Then:
%     \begin{enumerate}
%         \item $i(f,g) = l\left(g, \frac{S}{\langle f \rangle}\right)$,
%         \item $i(f,g)$ is finite if and only if $\sqrt{\langle f,g \rangle}=\mathfrak{n}$,
%         \item for every $h \in S$ we have  $i(f,gh)=i(f,g)+i(f,h)$,
%         \item $i(f,g) \geq \mult(f)\cdot \mult(g)$.
%     \end{enumerate}
% \end{lemma}

% \begin{proof}
    
% \end{proof}

%  In order to understand curve hypersurfaces, we present a version of the Weierstra{\ss} theorems over $S$. 

% \begin{definition}\label{def:hypersurface}
%     Let $f \in S$ and assume that $\mathfrak{n}$ is generated by $x$ and $y$:
%     \begin{enumerate}
%         \item We say that $f$ is \textbf{general with respect to $y$} if $f \in \mathfrak{n}$ but $f \notin \langle x \rangle$.  In this case, we define the order of $f$ with respect to $y$ to be $\ord_y(f)=i(f,x)$. 
%         \item We say that $f$ is a power series in $x$ if there exist some $u_i \in S$ that are either units or zero such that $f = \sum_{i=1}^\infty u_i x^i$. 
%     \end{enumerate}
% \end{definition}

% \begin{remark}
%  \textup{The second part of Definition~\ref{def:hypersurface} is the two dimensional version of Definition 3.5 in~\cite{svoray2025ade}. In that paper, the idea of power series with respect to a generator is used to preform an $ADE-$type classification for special hypersurface singularities. }
% \end{remark}

% \begin{lemma}\label{lem:weir_tech}
% \begin{enumerate}
%     \item Given $f \in \mathfrak{n}$ then we can find generators $x,y$ of $\mathfrak{n}$ such that $f$ is general with respect to $x$ and with respect to $y$.
%     \item Given some $x,y$ that generate $\mathfrak{n}$, then $f$ is general with respect to $y$ with $\ord_y(f)=r$ if and only if there exists a unit $u \in S$ such that $f-uy^r \in \langle x \rangle$.
% \end{enumerate}
% \end{lemma}

% \begin{proof}
%     For the first item, let $x_0$ and $y_0$ be generators of $\mathfrak{n}$. If $f \in \langle x_0 \rangle$ then we claim that there exists some $N \gg 0 $ such that $f \notin \langle x_0 + y_0^N\rangle$ (the result follows since $x_0 + y_0^N$ and $y_0+x_0^M$ generate $\mathfrak{n}$ for every $N,M>1$). 
% \end{proof}

% \begin{remark}
%     \textup{The second item in Lemma~\ref{lem:weir_tech} is the reason we call it "general with respect to $y$", even though the definition is given with based on $x$. It also relates our definition with the classical definition over power series, as defined in XXX.}
% \end{remark}




% \begin{proposition}\label{prop:Weierstrauss}
%     Assume that $x,y$ generate the maximal ideal of $S$ and let $f \in S$ be $y-$general with $r=\ord_y(f)$. Then:
%     \begin{enumerate}
%         \item (\textbf{Weierstrau{\ss} Division}) For every $g\in S$ there exists a unique $h \in S$ such that either $g=hf$ or 
%         \begin{equation*}
%             g-hf = u_{r-1}y^{r-1} + \cdots + u_0,
%         \end{equation*}
%         for some unique $u_0, \dots, u_{r-1} \in S$ that are power series in $x$. 
%         \item (\textbf{Weierstrau{\ss} Preperation}) We can find some unit $u \in S$ and some $u_1, \dots, u_r \in S$ that that are power series in $x$ such that
%         \begin{equation*}
%             f=u(y^{r}+u_{r-1}y^{r-1} + \cdots +u_0).
%         \end{equation*}
%     \end{enumerate}
%     In addition, in both cases, $u \in S$ and $u_1, \dots, u_r \in S$ are unique up to units. 
% \end{proposition}

% \begin{proof}
% \begin{enumerate}
%     \item 
%     \item By applying the first part to $g=y^r$ we can conclude that there exists some $h$ such that $y^r - hf = u_{r-1}y^{r-1}+\cdots +u_0$. Since $f$ is $y-$general then by Lemma~\ref{lem:weir_tech} we can write $f=uy^r+ax$ for some unit $u$ and some $a \in S$. Therefore we can conclude that $y^r -h(uy^r-ax)= u_{r-1}y^{r-1}+\cdots +u_0$. Since $y^r,f \in \mathfrak{m}$ and since $u_i$ is a power series in $x$ for every $i$, we can conclude that $u_i \in \langle x \rangle$. Therefore we can conclude that $y^r(1-hu) \in \langle x\rangle$, and since $x$ and $y$ form a regular sequence, we must have that $1-hu \in \langle x\rangle$. Therefore $h$ is a unit and we get that $f=h^{-1}(y^r-u_{r-1}y^{r-1}-\cdots -u_0)$, as desired. 
% \end{enumerate}    
% \end{proof}

% \begin{definition}
%     Assume that $x,y$  Assume that $x,y$ generate the maximal ideal of $S$ and let $f \in S$ be $y-$general and irreducible with $R= \frac{S}{\langle f \rangle}$. Then the \textbf{Puiseux uniformizer of $f$} is the unique uniformizer $\pi$ of $\overline{R}$ such that the image of $x$ under the map $\overline{\cdot } \colon S \twoheadrightarrow R \hookrightarrow \overline{R}$ is $\pi^{\ord_y(f)}$.  
% \end{definition}

% \begin{remark}
% \begin{enumerate}
%     \item \textup{Note that indeed the order of $\overline{x} \in \overline{R}$ is indeed $\ord_y(f)$ since.... }
%     \item \textup{The name Puiseux uniformizer comes from.... Parametrization and universal property. }
% \end{enumerate}
% \end{remark}

% \begin{proposition}\label{prop:prod_weier}
%     Assume that $x,y$  Assume that $x,y$ generate the maximal ideal of $S$ and let $f \in S$ be $y-$general and irreducible with $r=\ord_y(f)$. Let $\zeta$ be an $r-$th root of unity over $S$, and let $\xi \in \overline{R}$ satisfy $\xi^r = x$. Then we can find some unit $u$ and some elements $y_1, \dots, y_r \in R[\xi, \zeta]$ such that $f \in S \subset S[\xi, \zeta]$ factors as $f=u\prod_{i=1}^r (y-y_j)$ and $\overline{y_j}=u\left(\overline{\zeta}^j \overline{\xi}\right)^b \in \overline{R}$ for every $j$, where $\overline{y}=u\pi^b$ and $\pi$ is the Puiseux uniformizer of $f$. 
% \end{proposition}

% \begin{proof}
    
% \end{proof}

% \begin{remark}
%     \textup{Proposition~\ref{prop:prod_weier} is an analgue of Proposition 3.4 in~\cite{greuel2007introduction}, which allows us to factor $f$ up to roots of $x$.}
% \end{remark}

% We now use the "monic form" of $f$ given by the Weierstrau{\ss} preperation theorem in order to define an analogue of derivative of $f$ we respect to $y$. 

% \begin{definition}\label{def:Weierpar}
%     Assume that $x,y$ generate the maximal ideal of $S$ and let $f \in S$ be $y-$general. We set the \textbf{Weierstrau{\ss} derivative} of $f$ with respect to $y$ to be $\parwy(f)=u(ry^{r-1}+ (r-1)u_{r-1}y^{r-2}+ \cdots +u_1)$ where  $f=u(y^{r}+u_{r-1}y^{r-1} + \cdots +u_0)$ as in Proposition~\ref{prop:Weierstrauss}. 
% \end{definition}

% \begin{remark}
%     \textup{Note that by Proposition~\ref{def:Weierpar} we can conclude that $\parwy(f)$, as defined in Definition~\ref{def:Weierpar}, is well defined up to a unit. Therefore, it is well defined to look at the intersection multiplicity of $\parwy(f)$ with different elements of $S$. }
% \end{remark}

% \begin{lemma}
% Assume that $x,y$ generate the maximal ideal of $S$ and let $f \in S$ be $y-$general. Then:
%     \begin{enumerate}
%         \item $i(f,x) \leq i(\parwy(f),x)+1$.
%         \item $\parwy(f)$ is $y-$general if and only if $r \notin  \langle x \rangle$, which is equivalent to having equality in the previous item. 
%     \end{enumerate}
% \end{lemma}

% \begin{proof}
    
% \end{proof}

% \begin{remark}
%     \textup{Lemma 2.1 in~\cite{barroso2018milnor}}
% \end{remark}

% \begin{theorem}[Teissier Lemma]
%     Assume that $x,y$ generate the maximal ideal of $S$ and let $f \in S$ be $y-$general. Then 
%     \begin{equation*}
%         i(f, \parwy(f)) \leq \mu(f) + \ord_y(f)-1,
%     \end{equation*}
%     with equality if and only if.....
% \end{theorem}

% \begin{proof}
    
% \end{proof}

% Therefore, inspired by 

% \begin{definition}
%     Assume that $x,y$ generate $\mathfrak{n}$ and assume that $f$ is both $x-$ and $y-$regular. Then we define that \textbf{Weierstrau{\ss}-Milnor number of $f$ with respect to $x$ and $y$} to be 
%     \begin{equation*}
%         \muxyw(f)=\length_S\left(\frac{S}{\langle \parwx(f) ,\parwy(f)\rangle}\right).
%     \end{equation*}
% \end{definition}

    
% \begin{proposition}
%     Assume that $x,y$ generate $\mathfrak{n}$ and assume that $f$ is both $x-$ and $y-$regular. Then $\muxyw(f) \geq \mu(f)$.  
% \end{proposition}

% \begin{proof}
    
% \end{proof}


%\begin{definition}
%Assume that $\text{char}(\kappa) \neq 2$ and let $\Lambda$ be a lift of $\kappa$ to $S$. 
%\begin{enumerate}
%    \item A morsification of a hypersurface $R$ is a deformation $\mathcal{R}=\frac{S[t]}{\langle f_t \rangle}$ such that there exists some open Zariski set $U \subset \Lambda$ such that $\mu((R_{t_0})_{\mathfrak{m}}) \in \{0,1\}$ for every maximal ideal $\mathfrak{m}$ of $R_{t_0}$. 
%    \item Given a morsification $\mathcal{R}$ and some $t_0 \in S$, the number of maximal ideals $\mathfrak{m}$ such that $\mu((R_{t_0})_{\mathfrak{m}})=1$ is denoted by $\# A_1(R_{t_0})$.
%\end{enumerate}
%\end{definition}

%\begin{theorem}
%    Assume that $\text{char}(\kappa) \neq 2$ and let $\Lambda$ be a lift of $\kappa$ to $S$. Then for every hypersurface $R=\frac{S}{\langle f \rangle}$ such that $f \in \mathfrak{n}^2$ there exists a morsification $\mathcal{R}$. In addition, there exists a open Zariski set $\Lambda \subset S$ such that $\mu(R)=\#A_1(R_{t_0})$ for every $t_0 \in U$. 
%\end{theorem}
%\begin{proof}
%    Define $f_t = f+ txy$, and observe that indeed $\mathcal{R}=\frac{S}{\langle f_t \rangle}$ is an unfolding. 
%\end{proof}


% \section{Unfoldings of Quotient Rings}

% Throughout this section, we focus on quotient local rings $(R, \mathfrak{m})$ of the form $R=\frac{S}{I}$, where $(S, \mathfrak{n})$ is a complete regular ring of dimension $n$ and $I \subset S$ is a proper ideal. 

% \begin{definition}
%     Let $R=\frac{S}{I}$ be a quotient local ring. An \textbf{unfolding of $R$} is a ring of the form $\mathcal{R}_t=\frac{S[t]}{I_t}$, where $I_t \subset S[t]$ is an ideal that satisfies $I_t \mod t = I$. Given some $t_0 \in S$, we denote $\mathcal{R}_{t_0} = \frac{S[t]}{I_t + \langle t-t_0 \rangle}$.  
% \end{definition}


% In this section we focus on quotient local rings, that is, local rings $(R, \mathfrak{m})$ such that $R=\frac{S}{I}$ for some complete regular local ring $(S, \mathfrak{n})$.

% \begin{definition}
%     Let $k$ be a non-negative integer. We say that a quotient local ring $\frac{S}{I}$ is  $k-$determined if for every $J \subset S$ with $\frac{S}{I + \mathfrak{n}^{k+1}} \cong \frac{S}{J + \mathfrak{n}^{k+1}}$ we have that $\frac{S}{I} \cong \frac{S}{J}$. We say that $R$ is finitely determined if it is $k-$determined for some $k \geq 0$. 
% \end{definition}




\bibliographystyle{alpha}
\bibliography{bib}


\end{document}
